\documentclass[11pt,a4paper]{book}
\usepackage[margin=1in]{geometry}
\usepackage[T1]{fontenc}
\usepackage[utf8]{inputenc}
\usepackage{lmodern}
\usepackage{amsmath,amssymb,mathtools,bm}
\usepackage{booktabs,tabularx,array,longtable}
\usepackage{hyperref,xcolor}
\usepackage[most]{tcolorbox}
\usepackage{enumitem,algorithm,algpseudocode,fancyhdr}
\hypersetup{colorlinks=true,linkcolor=blue!55!black,urlcolor=blue!55!black}
\newcommand{\E}{\mathbb{E}}
\newcommand{\Prob}{\mathbb{P}}
\newcommand{\D}{\mathcal{D}}
\newcommand{\R}{\mathbb{R}}
\newcommand{\States}{\mathcal{S}}
\newcommand{\Actions}{\mathcal{A}}
\newcommand{\TODO}[1]{\textcolor{red!70!black}{[To write: #1]}}
\newtcolorbox{intuition}{colback=blue!3,colframe=blue!40!black,title=Intuition,breakable}
\newtcolorbox{pitfall}{colback=orange!3,colframe=orange!45!black,title=Common Confusion,breakable}
\setlength{\headheight}{14pt}
\pagestyle{fancy}\fancyhf{}\fancyhead[LE,RO]{Deep RL Survey}\fancyfoot[C]{\thepage}
\title{Deep Reinforcement Learning\\\large From Core Ideas to Modern Algorithms}
\author{CS 8803 Learning Notes}
\date{Version 1.1 \quad October 2026}
\begin{document}
\frontmatter\maketitle
\chapter*{About This Document}
This is a growing set of study notes based mainly on the Georgia Tech CS 8803 Deep Reinforcement Learning lectures (Fall 2026). It is not yet a complete research survey. The language is intentionally simple. We focus on why each method exists, what it learns, how it is trained, and how it chooses actions. Unfinished sections are clearly marked.
\tableofcontents
\mainmatter
\chapter{A Map of Reinforcement Learning}
\section{The Shared Goal}
An agent observes state $s_t$, takes action $a_t$, receives reward $r_t$, and moves to the next state. A typical goal is to maximize expected return:
\begin{equation}J(\pi)=\E_{\tau\sim p_\pi(\tau)}\left[\sum_{t=0}^{T-1}\gamma^t r_t\right].\end{equation}
\section{Three Main Things a Model Can Learn}
\begin{description}[style=nextline]
\item[Policy, $\pi_\theta(a\mid s)$] Given a state, it tells us which action to take or how likely each action is.
\item[Value, $V^\pi(s)$ or $Q^\pi(s,a)$] It estimates future reward. With discrete actions, choose $\arg\max_a Q(s,a)$.
\item[Dynamics, $P(s'\mid s,a)$] It predicts the next state, not the best action by itself.
\end{description}
\begin{intuition}
Policy: What should I do?\\
Q-function: How good is this action?\\
Dynamics model: What happens after I do it?
\end{intuition}
\section{Common Methods}
\small
\begin{longtable}{p{0.20\textwidth}p{0.23\textwidth}p{0.47\textwidth}}
\toprule Method & What it learns & Main idea \\\midrule\endhead
REINFORCE / VPG & Policy & Improve a policy with sampled returns and gradients.\\
Actor-Critic & Actor and value/Q & A critic helps the actor learn.\\
GAE & Advantage estimate & Control bias and variance; not a full RL algorithm.\\
PPO & Policy and often value & Limit each policy update.\\
DQN & Q-function & Choose discrete actions using action values.\\
Double DQN & Q-function & Reduce max-related overestimation.\\
DDPG & Actor and Q & Off-policy learning for continuous actions.\\
TD3 & Actor and two Qs & Improve DDPG stability.\\
SAC & Stochastic actor and Qs & Use entropy to encourage exploration.\\
MPC / CEM & Dynamics and planning & Search over possible future actions.\\
\bottomrule
\end{longtable}
\normalsize
\section{Not Every Concept Is an RL Algorithm}
A replay buffer stores transitions. A target network stabilizes targets. GAE estimates advantages. CEM is an optimizer.
\section{Six Questions for Every Method}
Why does it exist? What are its networks? Where does data come from? What is the loss? Where does gradient flow? How are actions selected?
\chapter{RL Fundamentals：价值与学习信号}
\section{Markov Decision Process}
一个常见的 MDP 写为 $\mathcal M=(\Sset,\A,P,r,\gamma,\rho_0)$，其中 $P(s'\mid s,a)$ 表示状态转移，$r(s,a)$ 为期望即时奖励，$\gamma\in[0,1)$ 为折扣因子，$\rho_0$ 为初始状态分布。有限回合可采用相应的终止条件。
\begin{align}
 s_0 &\sim \rho_0, & a_t &\sim \pi(\cdot\mid s_t),\\
 s_{t+1}&\sim P(\cdot\mid s_t,a_t), & r_t&\sim p(\cdot\mid s_t,a_t,s_{t+1}).
\end{align}
\section{Return、State Value、Action Value}
从时间 $t$ 开始的折扣回报为
\begin{equation}G_t=\sum_{k=0}^{T-t-1}\gamma^k r_{t+k}.\end{equation}
相对一条固定策略 $\pi$，定义
\begin{align}
V^\pi(s)&=\E_\pi[G_t\mid s_t=s],\\
Q^\pi(s,a)&=\E_\pi[G_t\mid s_t=s,a_t=a],\\
A^\pi(s,a)&=Q^\pi(s,a)-V^\pi(s).
\end{align}
\begin{intuition}
$V^\pi(s)$：我目前处于这个状态，之后一直跟随 $\pi$，平均能拿到多少回报？
$Q^\pi(s,a)$：我现在\emph{先做指定动作 $a$}，之后再跟随 $\pi$，回报如何？
$A^\pi(s,a)$：这个动作是否比该状态下的策略平均水平更好？
\end{intuition}
\section{Bellman Expectation vs. Bellman Optimality}
对已给定的策略 $\pi$，Bellman Expectation Equation 为
\begin{align}\label{eq:bellman-v}
V^\pi(s)&=\E_{a\sim\pi,\,s'\sim P}\big[r(s,a)+\gamma V^\pi(s')\big],\\
Q^\pi(s,a)&=\E_{s'\sim P}\left[r(s,a)+\gamma\E_{a'\sim\pi(\cdot\mid s')}Q^\pi(s',a')\right].
\end{align}
对最优策略，对应 Bellman Optimality Equation：
\begin{equation}\label{eq:bellman-opt}
Q^*(s,a)=\E_{s'\sim P}\left[r(s,a)+\gamma\max_{a'}Q^*(s',a')\right].
\end{equation}
这里的区别至关重要：\textbf{Policy Evaluation} 计算固定策略的好坏；\textbf{Policy Improvement} 借助 Value/Q 寻找更优动作；最优 Bellman 备份在后续动作上使用 $\max$。
\section{Monte Carlo 与 Temporal Difference}
不知道真正的 $V^\pi$ 时，可以从轨迹构造学习信号。Monte Carlo 使用真实采样的未来回报：
\begin{equation}y_t^{\mathrm{MC}}=G_t.\end{equation}
一步 TD 使用即时奖励和\emph{当前估计的未来价值}：
\begin{equation}\label{eq:td}
y_t^{\mathrm{TD}}=r_t+\gamma(1-d_t)V_{\bar\phi}(s_{t+1}), \qquad
\delta_t=y_t^{\mathrm{TD}}-V_\phi(s_t).
\end{equation}
$d_t$ 在此表示真正终止状态指示；时间截断与真正终止需按照实现约定分别处理。用于回归的典型 Critic Loss 为
\begin{equation}\mathcal L_V(\phi)=\E_{(s_t,r_t,s_{t+1})\sim\D}\!\left[(V_\phi(s_t)-\operatorname{stopgrad}(y_t^{\mathrm{TD}}))^2\right].\end{equation}
\begin{pitfall}
TD Target 并不是“真正的价值 Ground Truth”，而是通过 Bootstrap 构造的学习目标。训练 Value Network 不是把动作与网络输出做均方误差，而是把价值预测与\emph{目标回报}做回归。Target 中的未来价值一般不反向传播。
\end{pitfall}
\section{一个具体数字示例}
假设 $r_t=2$, $\gamma=0.9$, $V_{\bar\phi}(s_{t+1})=5$, $V_\phi(s_t)=8$，且不是终止状态，则
\begin{equation}y_t^{\mathrm{TD}}=2+0.9\cdot5=6.5,\qquad\delta_t=6.5-8=-1.5.\end{equation}
在其他因素固定的情况下，Value Network 对该样本的预测应向 $6.5$ 靠近。这里 $6.5$ 仍是\emph{估计出的 target}，不是已知最优值。
\section{Policy Gradient 与 Advantage 的关系}
一个常见的 Policy Gradient 形式是
\begin{equation}\label{eq:pg}
\nabla_\theta J(\theta)=\E_\pi\left[\sum_t \nabla_\theta\log\pi_\theta(a_t\mid s_t)A^\pi(s_t,a_t)\right],
\end{equation}
其中具体时序权重与采样分布需和所用目标函数的约定一致。优势为正，倾向提高采样动作的概率；优势为负，倾向降低概率。Critic 则通过 Value/Q 的回归 Loss 学习，而非直接使用 Actor 的目标。
\section{GAE 预告：不是单独的网络}
Generalized Advantage Estimation 在给定当前轨迹和固定 Value 估计的情况下，结合多个时间跨度的 TD Residual：
\begin{equation}\hat A_t^{\mathrm{GAE}(\gamma,\lambda)}=\sum_{l=0}^{T-t-1}(\gamma\lambda)^l\delta_{t+l},\end{equation}
上式假设在轨迹结束处作合适的截断，实际计算还需要正确的 terminal mask。$\lambda$ 控制更多使用短期 bootstrap 或更长期的采样回报，从而影响偏差和方差。\textbf{GAE 通常在一批轨迹收集完之后先计算；PPO 对同一批数据做多次 minibatch 更新时，计算好的 Advantage 往往保持固定。}新一批轨迹则重新计算。
\section{本章自测}
\begin{enumerate}
\item 为什么 $Q^\pi(s,a)$ 比 $V^\pi(s)$ 多一个 Action 输入？
\item TD Target $r+\gamma V(s')$ 中的 $V(s')$ 不准确时，会发生什么？
\item 为什么 GAE 是估计器而不是新的 Policy 或 Dynamics Network？
\end{enumerate}
\chapter{Policy Optimization: From REINFORCE to PPO}
\section{Learning Goals and Algorithm Roadmap}
This chapter explains the progression
\begin{center}
REINFORCE $\rightarrow$ Reward-to-Go $\rightarrow$ Baselines
$\rightarrow$ Actor-Critic $\rightarrow$ GAE $\rightarrow$ PPO.
\end{center}
Each step addresses a specific problem: noisy gradients, credit assignment, value estimation, or overly large updates from reused rollout data. These are related ideas, not necessarily one historical chain of mandatory components.

\section{Why Learn a Policy Directly?}
A value-based discrete-action agent can select $\arg\max_a Q(s,a)$. This is difficult when actions are continuous and a global action maximization must be solved at each step. A policy $\pi_\theta(a\mid s)$ instead outputs an action distribution directly; a deterministic policy $\mu_\theta(s)$ outputs an action. Stochastic policies naturally support exploration.

\section{The Trajectory Probability and Log-Derivative Trick}
For a finite-horizon MDP with dynamics independent of $\theta$,
\begin{equation}
p_\theta(\tau)=\rho_0(s_0)\prod_{t=0}^{T-1}
\pi_\theta(a_t\mid s_t)P(s_{t+1}\mid s_t,a_t).
\end{equation}
For a sampled trajectory $\tau$ define $R(\tau)=\sum_{t=0}^{T-1}r_t$ in the undiscounted case.
\begin{align}
J(\theta)&=\E_{\tau\sim p_\theta}[R(\tau)],\\
\nabla_\theta p_\theta(\tau)&=
p_\theta(\tau)\nabla_\theta\log p_\theta(\tau),\\
\nabla_\theta\log p_\theta(\tau)&=
\sum_{t=0}^{T-1}\nabla_\theta\log\pi_\theta(a_t\mid s_t),\\
\nabla_\theta J(\theta)&=
\E_\tau\left[R(\tau)\sum_t\nabla_\theta\log\pi_\theta(a_t\mid s_t)\right].
\end{align}
The initial-state and environment transition terms disappear from the log-probability gradient because they do not directly depend on $\theta$ for a fixed sampled trajectory. The environment still affects which trajectories and rewards are observed; its mathematical derivatives are not required.

\section{REINFORCE and Monte Carlo Estimation}
Sample $N$ trajectories under the current policy. A Monte Carlo policy-gradient estimator is
\begin{equation}
\hat g=\frac1N\sum_{i=1}^{N}R(\tau^{(i)})
\sum_{t=0}^{T-1}\nabla_\theta
\log\pi_\theta(a_t^{(i)}\mid s_t^{(i)}).
\end{equation}
With gradient descent, use the negative of the sample log-probabilities weighted by sampled returns. Returns are treated as constants in the actor computation graph. REINFORCE is on-policy and model-free; a critic is not required. The estimator can have high variance.

\section{Reward-to-Go and Causality}
The total return includes rewards before action $a_t$, which that action cannot influence. The policy score satisfies
\begin{equation}
\E_{a_t\sim\pi_\theta(\cdot\mid s_t)}
[\nabla_\theta\log\pi_\theta(a_t\mid s_t)]=0.
\end{equation}
Conditioning on the history before $a_t$ shows that adding past rewards to the weight has zero expected gradient contribution. We can use future rewards only:
\begin{equation}
G_t=\sum_{k=t}^{T-1}\gamma^{k-t}r_k.
\end{equation}
For the undiscounted objective, a reward-to-go estimator is $\sum_t\nabla\log\pi_\theta(a_t\mid s_t)G_t$. For the standard discounted objective $J=\E[\sum_t\gamma^tr_t]$, the corresponding expression has an additional $\gamma^t$ outside this product; algorithms may use other objective conventions.
Example: rewards $(1,2,10)$ with $\gamma=1$ yield full-return weights $(13,13,13)$ versus reward-to-go weights $(13,12,10)$. With rewards $(5,-2,8)$ and $\gamma=0.9$, reward-to-go values are $(9.68,5.2,8)$.
\begin{intuition}
An action should not receive gradient noise from rewards that occurred before the action.
\end{intuition}

\section{Baselines and the Advantage}
For any action-independent state baseline $b(s_t)$, conditioning on $s_t$ gives
\begin{align}
\E_{a_t\sim\pi_\theta}[
\nabla_\theta\log\pi_\theta(a_t\mid s_t)b(s_t)]
&=b(s_t)\nabla_\theta\sum_a\pi_\theta(a\mid s_t)=0.
\end{align}
Therefore replacing $G_t$ by $G_t-b(s_t)$ preserves the expected policy-gradient contribution, under the usual sampling and differentiability assumptions. A poor baseline need not reduce variance. A natural choice is $b(s)=V^\pi(s)$:
\begin{equation}
A^\pi(s,a)=Q^\pi(s,a)-V^\pi(s),\qquad
\hat A_t^{MC}=G_t-V_\phi(s_t).
\end{equation}
For example, at a state with $Q(L)=100$, $Q(R)=110$, and an equal-probability policy, $V(s)=105$ and advantages are $-5,+5$. The advantage measures relative action quality; the sampled return $G_t$ is still noisy.

\section{Actor-Critic Architecture and One-Step TD Advantage}
An independent actor $\pi_\theta(a\mid s)$ outputs action probabilities. A state-value critic $V_\phi(s)$ outputs a scalar. For nonterminal transitions:
\begin{align}
y_t^{TD}&=r_t+\gamma V_\phi(s_{t+1}),\\
\delta_t&=y_t^{TD}-V_\phi(s_t),\\
L_V(\phi)&=\E[(V_\phi(s_t)-\operatorname{stopgrad}(y_t^{TD}))^2],\\
L_{\mathrm{actor}}(\theta)&=-\E[
\log\pi_\theta(a_t\mid s_t)\operatorname{stopgrad}(\delta_t)].
\end{align}
With an exact $V^\pi$,
$\E[\delta_t\mid s_t,a_t]=A^\pi(s_t,a_t)$.
With an approximate critic this may be biased. Example: $r=2$, $\gamma=0.9$, $V(s)=4$, $V(s')=6$; $y^{TD}=7.4$, $\delta=3.4$. The critic moves toward $7.4$; the actor tends to increase the sampled action probability because the advantage estimate is positive. Genuine terminal transitions have zero next-state bootstrap contribution.

\section{Multi-Step Returns and Generalized Advantage Estimation}
A one-step advantage relies strongly on the critic's next-state prediction; a full Monte Carlo advantage can have high variance. For an episode with at least $n$ remaining steps,
\begin{align}
\hat Q_t^{(n)}&=\sum_{k=0}^{n-1}\gamma^kr_{t+k}
+\gamma^n V_{\mathrm{old}}(s_{t+n}),\\
\hat A_t^{(n)}&=\hat Q_t^{(n)}-V_{\mathrm{old}}(s_t)
=\sum_{l=0}^{n-1}\gamma^l\delta_{t+l}.
\end{align}
GAE combines TD residuals:
\begin{equation}
\hat A_t^{GAE(\gamma,\lambda)}
=\sum_{l=0}^{T-t-1}(\gamma\lambda)^l\delta_{t+l}.
\end{equation}
Use proper terminal masking and bootstrap at an unfinished rollout boundary. When $\lambda=0$, GAE is one-step TD advantage. When $\lambda=1$ and the episode truly ends, it equals $G_t-V_{\mathrm{old}}(s_t)$. Increasing $\lambda$ usually lowers bootstrap bias but raises variance.
For $(\delta_0,\delta_1,\delta_2)=(2,4,-1)$, $\gamma=0.9$, $\lambda=0.5$:
\begin{equation}
\hat A_0=2+0.45(4)+0.45^2(-1)=3.5975.
\end{equation}

\section{GAE and TD($\lambda$)}
The forward-view TD($\lambda$) target is the lambda-return $G_t^\lambda$. When computed using the same fixed rollout value estimates,
\begin{equation}
G_t^\lambda=V_{\mathrm{old}}(s_t)+\hat A_t^{GAE}.
\end{equation}
A finite-episode mixture of $n$-step returns, with $N$ steps remaining, is
\begin{equation}
G_t^\lambda=(1-\lambda)\sum_{n=1}^{N-1}\lambda^{n-1}G_t^{(n)}
+\lambda^{N-1}G_t^{(N)}.
\end{equation}
The forward-view lambda-return is a \emph{critic target}; GAE is an \emph{actor advantage estimate}. Classical backward-view TD($\lambda$) additionally involves eligibility traces.
For the previous example, if $V_{\mathrm{old}}(s_0)=5$, then $G_0^\lambda=8.5975$.

\section{On-Policy Data Reuse and Importance Ratios}
PPO gathers a batch with behavior policy $\pi_{\theta_{\mathrm{old}}}$ and saves the sampled actions and their old log-probabilities. During optimization, the current policy $\pi_\theta$ changes, but the stored actions do not.
The importance ratio for each sampled action is
\begin{equation}
\rho_t(\theta)=\frac{\pi_\theta(a_t\mid s_t)}
{\pi_{\theta_{\mathrm{old}}}(a_t\mid s_t)}
=\exp(\log\pi_\theta(a_t\mid s_t)
-\log\pi_{\theta_{\mathrm{old}}}(a_t\mid s_t)).
\end{equation}
The importance sampling identity changes expectations between two action distributions, provided their support condition holds. The PPO surrogate $\E[\rho_t\hat A_t]$ uses the \emph{old state distribution}, so it is not an exact arbitrary-distance correction of full new-policy trajectory returns.
\begin{pitfall}
At the next PPO minibatch update, recompute the \emph{new log-probability of the stored action}. Do not sample a replacement action from the current policy. Keep old log-probabilities, advantages, and returns fixed for this rollout.
\end{pitfall}

\section{PPO Clipped Objective}
PPO maximizes
\begin{equation}
L^{CLIP}(\theta)=\E_t\left[
\min\left(
\rho_t\hat A_t,\;
\operatorname{clip}(\rho_t,1-\epsilon,1+\epsilon)\hat A_t
\right)\right].
\end{equation}
The minimum clips \emph{excessive improvement}, not every ratio beyond the interval:
\begin{itemize}
\item If $\hat A_t>0$, the upper side $\rho_t>1+\epsilon$ is capped.
\item If $\hat A_t<0$, the lower side $\rho_t<1-\epsilon$ is capped.
\item If a good action becomes much less probable, or a bad action much more probable, the objective retains a corrective gradient (from the un-clipped term).
\end{itemize}
For $\epsilon=0.2$, $\rho=1.5$, $\hat A=-4$: the terms are $-6$ and $-4.8$. The PPO objective is $\min(-6,-4.8)=-6$, not $-4$ or $-4.8$.
A naive objective $\operatorname{clip}(\rho,1-\epsilon,1+\epsilon)\hat A$ incorrectly removes corrective gradients on both sides.
Clipping is a practical conservative surrogate, \emph{not} a hard guarantee that all ratios or policy KL values are bounded.

\section{PPO Complete Data and Gradient Flow}
\begin{enumerate}
\item Collect a rollout with $\pi_{\theta_{\mathrm{old}}}$. Save states, sampled actions, rewards, next-state information, old action log-probabilities, and rollout-time value predictions.
\item With old values, compute TD residuals, GAE advantages, and $\hat R_t=\hat A_t+V_{\mathrm{old}}(s_t)$. Freeze these batch targets.
\item Repeat several minibatch epochs. Feed \emph{stored} states and actions into the \emph{current} actor and recompute $\log\pi_\theta(a_t\mid s_t)$.
\item Compute ratios and clipped policy objective, then update actor parameters $\theta$.
\item Predict $V_\phi(s_t)$ using the current critic, regress toward fixed $\hat R_t$, and update critic parameters $\phi$.
\item After optimization, collect a new rollout under the updated policy, reset the reference old policy, and recompute targets.
\end{enumerate}
For independent actor and critic networks, detached advantages make actor-loss gradient flow only to $\theta$, and the value loss updates $\phi$. If actor and critic share a backbone, both losses may update shared features.
\subsection{Minimal Code Data Flow}
\begin{verbatim}
# after rollout: all tensors [B], except states [B, obs_dim]
with torch.no_grad():
    advantages = compute_gae(rewards, old_values, boundaries)
    returns = advantages + old_values
    old_log_probs = saved_behavior_log_probs
for epoch in range(n_epochs):
    for batch in minibatches:
        new_log_probs = actor(states).log_prob(stored_actions)
        ratio = exp(new_log_probs - old_log_probs)
        actor_loss = -min(ratio * advantages,
            clamp(ratio, 1-eps, 1+eps) * advantages).mean()
        value_loss = mse(critic(states), returns)
        # Step independent actor and critic optimizers.
\end{verbatim}
This is schematic: indices, shapes, rollout boundary bootstrap, optional entropy term, and optional KL monitoring must be handled in a working implementation.
\section{Review Questions}
\begin{enumerate}
\item Why can trajectory policy gradients ignore the derivatives of environment dynamics?
\item Why can reward-to-go remove earlier rewards without changing the expected gradient?
\item Why does a state baseline preserve the expected policy gradient?
\item When is TD residual an unbiased conditional estimator of advantage?
\item How does GAE interpolate between one-step TD and Monte Carlo?
\item Why is $\hat R_t$ kept fixed over the minibatch optimization epochs?
\item For $\rho=1.5$ and $\hat A=-4$ with $\epsilon=0.2$, what does the PPO minimum choose?
\item Why must the next optimization epoch reuse the old sampled action?
\end{enumerate}
\section{Learning Status and Source Notes}
First study pass completed. Revisit the distinction between \emph{stored action} and \emph{new action probability}, and the negative-advantage side of PPO clipping.
Course basis: supplied CS 8803 Fall 2026 Lecture 2 (Policy Gradients) and Lecture 3 (Actor-Critic Algorithms): likelihood-ratio gradients, reward-to-go, baselines, on/off-policy importance sampling, critics, multi-step returns, and GAE. Detailed PPO clipped-objective walkthrough and implementation are supplementary study explanation; the above is not a verbatim slide transcription.

\chapter{Value-Based RL: Q-Learning and DQN}
\section{Learning Goals and Roadmap}
Q-Learning $\rightarrow$ neural Q-functions $\rightarrow$ DQN
$\rightarrow$ Double DQN. We study action selection, data, targets,
losses, gradient flow, and limitations.

\section{Why Do We Need It?}
The objective is still to maximize expected discounted return.
A state value $V^\pi(s)$ evaluates a state, but alone does not rank
actions without additional knowledge of rewards and dynamics.
An action value evaluates a particular first action:
\begin{equation}
Q^\pi(s,a)=\E_\pi\left[\sum_{k=0}^{\infty}\gamma^k r_{t+k}
\mid s_t=s,a_t=a\right].
\end{equation}
After that first action, the agent follows $\pi$. The optimal
action value $Q^*(s,a)$ assumes optimal subsequent behavior.
For finite discrete actions, a greedy policy selects
$\arg\max_a Q(s,a)$. This is the best action according to the
current estimate, not necessarily the truly optimal action.

\section{Tabular Q-Learning}
Store one estimate for each state-action pair. For a sampled
transition $(s,a,r,s',d)$, with $d=1$ for genuine termination,
\begin{align}
y &= r+\gamma(1-d)\max_{a'}Q(s',a'),\\
\delta &= y-Q(s,a),\\
Q(s,a)&\leftarrow Q(s,a)+\alpha\delta.
\end{align}
The Bellman optimality equation motivates this sampled target:
\begin{equation}
Q^*(s,a)=\E_{r,s'\mid s,a}
[r+\gamma(1-d)\max_{a'}Q^*(s',a')].
\end{equation}
Example: $Q(s,a)=3$, $r=2$, $\gamma=0.9$, and next-state
values $(4,6)$. Then $y=7.4$, $\delta=4.4$, and with
$\alpha=0.1$ the updated entry is $3.44$.
The maximum expresses optimal subsequent behavior; averaging
action values would correspond to a different subsequent policy.

\subsection{Exploration and Off-Policy Learning}
An epsilon-greedy behavior policy chooses a random action with
probability $\epsilon$ and a greedy action otherwise.
The update still uses the greedy next-state value even if the
next action actually taken is exploratory. Thus data collection
and the target policy can differ: Q-learning is off-policy.
If next-state values are $(4,6)$ and exploration chooses the first
action, the previous transition's bootstrap value remains $6$.

\section{Models and Inputs/Outputs}
A neural Q-function replaces a table when states are high-dimensional
or too numerous. A common discrete-action architecture maps
$s\in\R^{n_s}$ to a vector in $\R^{n_a}$, with one Q-value per action.
These are return estimates, not probabilities; no softmax is needed.
For a batch, states have shape $[B,n_s]$ and outputs $[B,n_a]$.
For each transition, select the value of the action actually executed,
not the largest current-state output.
Example: outputs $(2,3)$ and stored action index $0$ give prediction $2$.
Shared parameters enable generalization but also make updates affect
other states and actions.

\section{DQN: Replay and Target Networks}
DQN combines a neural Q-function with experience replay and a slowly
updated target network.
\begin{description}[style=nextline]
\item[Online network $Q_\theta$] Selects interaction actions and produces
current predictions. Its parameters are updated by gradient descent.
\item[Target network $Q_{\bar\theta}$] Produces bootstrap estimates.
It is initialized from the online network and periodically refreshed
by $\bar\theta\leftarrow\theta$ in the hard-update version studied here.
\item[Replay buffer] Stores observed transitions $(s,a,r,s',d)$.
Random minibatches reduce temporal correlation and reuse experience.
Stored transitions remain fixed; targets are computed when sampled.
\end{description}
\begin{align}
y_i&=r_i+\gamma(1-d_i)\max_{a'}Q_{\bar\theta}(s'_i,a'),\\
L(\theta)&=\frac1B\sum_i
\left[Q_\theta(s_i,a_i)-\operatorname{stopgrad}(y_i)\right]^2.
\end{align}
Squared error illustrates the principle; Huber loss is also common.
Gradient flows through the current prediction only.
Detaching a target prevents differentiation through it in that update,
but does not make its numerical value permanently fixed.

\subsection{Complete Training Flow}
\begin{enumerate}
\item Use the online network with epsilon-greedy exploration to act.
\item Observe the transition and store it in replay.
\item Randomly sample a minibatch, possibly containing old experience.
\item Compute executed-action predictions with the online network.
\item Compute next-state targets with the target network, without gradients.
\item Backpropagate the loss and update the online network.
\item At the chosen interval, copy online parameters into the target network.
\end{enumerate}
Interaction and optimization need not use the same transition.
Updating Q-values implicitly changes the greedy policy: values $(5,3)$
choose left, while later values $(5,7)$ choose right, without an actor update.

\subsection{Minimal Code Data Flow}
\begin{verbatim}
# states, next_states: [B, obs_dim]
# actions: [B, 1], integer indices
# rewards, terminated: [B, 1]
q_pred = online_net(states).gather(1, actions)
with torch.no_grad():
    next_best = target_net(next_states).max(
        dim=1, keepdim=True).values
    target = rewards + gamma * (
        1 - terminated.float()) * next_best
loss = ((q_pred - target) ** 2).mean()
optimizer.zero_grad()
loss.backward()
optimizer.step()
# Periodically: target_net.load_state_dict(online_net.state_dict())
\end{verbatim}
This is schematic code. Handle environment resets and truncations
correctly: a time-limit truncation is not automatically a genuine
terminal state, and bootstrapping may require its final observation.

\begin{pitfall}
For ordinary DQN, a fixed transition has a fixed target between
hard target-network updates, assuming deterministic network outputs.
When the target network is refreshed, the target can change even
though the observed reward has not changed.
\end{pitfall}

\section{Double DQN: Separate Selection and Evaluation}
Taking a maximum over noisy estimates can cause overestimation.
If both true action values are $5$, estimates alternating between
$(4,6)$ and $(6,4)$ are individually unbiased, but their maximum
is always $6$.
Ordinary DQN uses the target network both to select the largest
next-state action value and to evaluate that action.
Double DQN instead uses:
\begin{align}
a_i^*&=\arg\max_{a'}Q_\theta(s'_i,a'),\\
y_i^{Double}&=r_i+\gamma(1-d_i)Q_{\bar\theta}(s'_i,a_i^*).
\end{align}
The existing two networks and their update schedule can remain the same;
the change is the construction of the target.
For online values $(8,3)$ and target values $(4,7)$, DQN bootstraps
with $7$, while Double DQN selects the first action online and evaluates
it as $4$ using the target network.
Separating selection and evaluation reduces max-related overestimation,
but does not guarantee removal of all bias; network errors remain correlated.
If both networks are identical, their targets coincide.

\begin{verbatim}
with torch.no_grad():
    next_actions = online_net(next_states).argmax(
        dim=1, keepdim=True)                 # [B, 1]
    next_q = target_net(next_states).gather(
        1, next_actions)                     # [B, 1]
    target = rewards + gamma * (
        1 - terminated.float()) * next_q
\end{verbatim}
The prediction and loss are unchanged.
Stored actions index actual past behavior at the current state.
Next actions are freshly selected to construct next-state targets.
Both operations above occur without gradients.

\begin{pitfall}
Double DQN targets can change even while target-network parameters
are fixed, because the online-selected action can change.
If target values remain $(4,7)$ while online values change from
$(8,3)$ to $(2,9)$, the bootstrap value changes from $4$ to $7$.
This selection concerns the training target. Environment interaction
still ordinarily uses online epsilon-greedy action selection.
\end{pitfall}

\section{Links to Other Methods and Limitations}
Yesterday's stochastic actor--state-value critic learns a policy
distribution and a scalar $V(s)$. The critic helps train the actor
through advantage estimates. DQN learns action values and derives
a greedy policy directly. Actor-critic critics may also be Q-functions,
and actors need not always be stochastic.
PPO reuses a recent rollout with fixed batch advantages and return
targets over optimization epochs; DQN reuses replay transitions
and recomputes bootstrap targets when sampled.

Standard DQN enumerates finite actions. Continuous actions would require
solving $\arg\max_a Q(s,a)$ over a continuous space at each decision.
The next chapter introduces actors that directly output continuous
actions and learn to increase their critic's estimated Q-value.
For a deterministic actor, this motivates
$J(\theta)=\E_s[Q_\phi(s,\mu_\theta(s))]$.

Neural approximation, bootstrapping, and off-policy learning together
form the deadly triad: their combination can cause instability or
divergence. Replay, target networks, and Double DQN improve practical
learning but do not provide general convergence guarantees for
nonlinear neural Q-learning. Exploration and data coverage also matter.

\section{Review Questions}
\begin{enumerate}
\item Why can a Q-function directly rank actions while a state value alone cannot?
\item Why does Q-learning take a next-action maximum instead of an average?
\item Why can an exploratory transition train a greedy target policy?
\item Which current-state output enters the loss, and why?
\item How does no-gradient target computation differ from freezing its values?
\item Why can the same replay transition receive a new DQN target?
\item What are the separate roles of replay and the target network?
\item In Double DQN, which network selects and which evaluates the next action?
\item Can a Double DQN target change before the target network is refreshed?
\item Why is continuous-action maximization harder, and how might an actor help?
\end{enumerate}

\section{Learning Status and Source Notes}
First study pass completed on October 13, 2026.
Clarified distinctions: executed action versus target action;
ordinary DQN's two networks versus Double DQN's selection/evaluation split;
fixed transitions versus recomputed targets; and implicit greedy-policy
improvement versus explicit actor learning.
Course basis: supplied Lectures 4 and 5 on value-based methods and
deep Q-functions, supplemented by our guided explanations and examples.
This chapter is a first-pass study note, not an exhaustive account of
all lecture topics or a verbatim slide transcription.

\chapter{Continuous Control: DDPG, TD3, SAC}
\section{Learning Goals and Roadmap}
DDPG learns an actor for continuous-action Q-learning.
TD3 improves critic targets and actor-update stability.
SAC learns a stochastic policy under an entropy-regularized objective.
We distinguish stored actions, current-policy actions, target actions,
and their gradient paths.

\section{Why Continuous Actions Need an Actor}
Discrete DQN outputs one value per action and selects an argmax.
Continuous actions cannot be fully enumerated. A network $Q_\phi(s,a)$
can evaluate an action, but solving $\arg\max_a Q_\phi(s,a)$ at every
decision can be expensive. DDPG learns an approximate maximizer
$\mu_\theta(s)$ that directly outputs a continuous action.
This does not guarantee global maximization.
All three methods are model-free and off-policy and use replay.

\section{DDPG: Networks, Data, and Targets}
DDPG means Deep Deterministic Policy Gradient.
The online actor maps states $[B,n_s]$ to actions $[B,n_a]$.
The online Q-critic takes states and actions and outputs $[B,1]$.
A critic is a role: it can be $V(s)$ or $Q(s,a)$, with different inputs.
DDPG maintains target copies of both networks, for four networks total.

Interaction typically uses
\begin{equation}
a_t=\operatorname{clip}(\mu_\theta(s_t)+\epsilon_t,a_{\min},a_{\max}).
\end{equation}
Replay stores this actual noisy, clipped action with its observed
reward, next state, and genuine-termination flag.
For sampled $(s,a,r,s',d)$:
\begin{align}
a'&=\mu_{\bar\theta}(s'),\\
y&=r+\gamma(1-d)Q_{\bar\phi}(s',a'),\\
L_Q(\phi)&=\E[(Q_\phi(s,a)-\operatorname{stopgrad}(y))^2].
\end{align}
The current prediction uses the stored action. Replacing it with
a new actor action would mismatch the observed reward and next state.
The target is partly observation and partly prediction, not an
entirely ground-truth return.

\subsection{Actor Loss and Gradient Source}
For replay states, propose fresh actor actions:
\begin{align}
J_\mu(\theta)&=\E_{s\sim\D}[Q_\phi(s,\mu_\theta(s))],\\
L_\mu(\theta)&=-J_\mu(\theta),\\
\nabla_\theta J_\mu
&=\E_s\left[
\left.\nabla_a Q_\phi(s,a)\right|_{a=\mu_\theta(s)}
\nabla_\theta\mu_\theta(s)\right].
\end{align}
This is the practical replay-based objective.
Fix critic parameters during actor optimization, but retain derivatives
with respect to its action input. No environment derivative is needed.
A critic $Q(s,a)=10-(a-0.7)^2$ has action derivative $0.8$ at $a=0.3$,
guiding a local increase in the actor's output. Wrong learned slopes
can mislead the actor.

\begin{pitfall}
Freezing critic parameters preserves the action-input gradient.
Detaching an actor action preserves its value but cuts its connection
to actor parameters. Detach does not remove the action from the input
or convert a Q-critic into a V-critic.
\end{pitfall}

\begin{verbatim}
with torch.no_grad():
    next_a = target_actor(next_states)
    y = rewards + gamma * (1 - terminated.float()) * \
        target_critic(next_states, next_a)
critic_loss = (critic(states, stored_actions) - y).square().mean()
critic_optimizer.zero_grad()
critic_loss.backward()
critic_optimizer.step()

for p in critic.parameters():
    p.requires_grad_(False)
actor_optimizer.zero_grad()
actor_loss = -critic(states, actor(states)).mean()
actor_loss.backward()
actor_optimizer.step()
for p in critic.parameters():
    p.requires_grad_(True)
\end{verbatim}

\subsection{Target Actor and Soft Updates}
The online actor learns through the online critic.
The target actor only supplies next-state actions for critic labels.
It has no separate actor loss. Slowly follow online parameters:
\begin{equation}
\bar\theta\leftarrow(1-\tau)\bar\theta+\tau\theta,\qquad
\bar\phi\leftarrow(1-\tau)\bar\phi+\tau\phi.
\end{equation}
Even with a fixed target critic, changing its action input can change
the target value. A slow actor copy reduces that movement but lags
behind the newest policy.

\subsection{Connection to Yesterday's Actor-Critic}
A basic stochastic actor loss is
$-\E[\log\pi_\theta(a_{\rm stored}\mid s)\operatorname{stopgrad}(\hat A)]$.
It changes the probability of an observed action according to advantage.
DDPG changes a newly proposed action using the critic's action derivative.
For a Gaussian with fixed standard deviation,
$\partial\log\pi/\partial m=(a-m)/\sigma^2$:
positive advantage for $a=0.4$ with mean $0.3$ pulls the mean toward
the observed action. A V-critic action derivative is not required.
At a fixed state, $-V_\phi(s)$ contains no independent actor parameters
and cannot directly train that actor.

\section{TD3: Three Improvements}
TD3 means Twin Delayed Deep Deterministic Policy Gradient.
An actor maximizing an inaccurate critic can exploit spurious Q peaks,
and bootstrapping can propagate these errors.

\subsection{Clipped Double Q-Learning}
Train two online critics with separate parameters and random
initializations, and maintain their target copies:
\begin{align}
y&=r+\gamma(1-d)\min_{i=1,2}Q_{\bar\phi_i}(s',\tilde a'),\\
L_{Q_i}&=\E[(Q_{\phi_i}(s,a)-\operatorname{stopgrad}(y))^2].
\end{align}
Both regress toward the same label. For $r=2$, $\gamma=0.9$,
and target estimates $(12,8)$, the target is $9.2$.
Online predictions $(7,10)$ move upward and downward respectively.
The minimum reduces overestimation but can cause underestimation;
correlated errors remain possible.
Double DQN separates action selection from evaluation.
TD3 instead evaluates one actor-proposed action with two critics.

\subsection{Delayed Policy Updates}
Critics update every optimization step, but the actor less frequently.
For policy delay 2:
\begin{center}
\begin{tabular}{c l l l}
\toprule Step & Online critics & Online actor & All targets\\
\midrule
1 & Update & Skip & Skip\\
2 & Update & Update & Soft update\\
3 & Update & Skip & Skip\\
4 & Update & Update & Soft update\\
\bottomrule
\end{tabular}
\end{center}
Delay means frequency, not simply updating the critic before the actor.
Without delay both can update every step.
Extra critic updates allow evaluation to improve before guiding
the actor; this is not a guarantee of accuracy.
Standard TD3 trains the actor using online critic 1:
\begin{equation}
L_\mu=-\E[Q_{\phi_1}(s,\mu_\theta(s))].
\end{equation}
A loss using only target actor and target critic cannot train the
online actor: its parameters are absent from that computation graph.
An optimizer cannot create a missing gradient connection.

\subsection{Target Policy Smoothing}
\begin{align}
\epsilon&\sim\mathcal N(0,\sigma_{\rm target}^2I),\\
\tilde\epsilon&=\operatorname{clip}(\epsilon,-c,c),\\
\tilde a'&=\operatorname{clip}(
\mu_{\bar\theta}(s')+\tilde\epsilon,a_{\min},a_{\max}).
\end{align}
Both target critics evaluate the same perturbed action.
Repeated sampling smooths target learning near the proposed action,
reducing dependence on narrow erroneous Q peaks. Usually one
perturbation per transition is sampled, rather than explicitly
averaging an entire neighborhood.
Exploration noise affects actual environment actions;
target smoothing noise affects training labels only.

TD3 has six networks: one actor, two critics, and three target copies.
Collect online-actor noisy behavior, sample replay, construct
no-gradient smoothed twin-Q targets, and update both critics.
At the policy-delay interval, update the actor and soft-update all targets.
For a stored action $0.4$, a target next action $0.62$, and current actor
output $0.5$, the three actions serve critic regression, target
construction, and actor training respectively.

\section{SAC: Reward Plus Entropy}
SAC means Soft Actor-Critic. We study the common modern version:
one stochastic actor, two online Q-critics, two target Q-critics,
and replay, typically no target actor or separate V-network.
Earlier SAC formulations differ.
\begin{equation}
J(\pi)=\E_{\tau\sim\pi}\left[
\sum_t\gamma^t(r_t+\alpha\mathcal H(\pi(\cdot\mid s_t)))\right].
\end{equation}
Temperature $\alpha>0$ weights entropy relative to reward.
Entropy encourages maintaining alternative actions, not unconditional
random behavior. A stochastic actor does not imply on-policy learning.

\subsection{Entropy, Density, and an Example}
\begin{equation}
\mathcal H(\pi(\cdot\mid s))=-\E_{a\sim\pi}[\log\pi(a\mid s)].
\end{equation}
For fixed discrete action values $(10,9)$, always choosing the first
gives value-plus-entropy objective $10$. Equal probabilities give
$9.5+\alpha\log2$, or $10.193$ at $\alpha=1$.
A large value gap can outweigh the entropy incentive.
This example illustrates a local policy-improvement objective.

Continuous $\pi(a\mid s)$ is a density, not a point probability.
A uniform density on $[0,0.1]$ equals $10$ and integrates to $1$;
a subinterval of width $0.01$ has probability $0.1$.
A Gaussian of standard deviation $0.1$ has peak density $3.99$,
with negative log-density approximately $-1.38$ there.
Differential entropy can be negative and depends on action coordinates.
A single action's density does not determine the whole policy's entropy.

\subsection{Actor Outputs and Reparameterization}
The actor usually outputs mean and log standard deviation,
both shaped $[B,n_a]$:
\begin{align}
\epsilon&\sim\mathcal N(0,I),\\
u&=m_\theta(s)+\sigma_\theta(s)\odot\epsilon,\\
a&=\tanh(u).
\end{align}
Noise is parameter-independent but actions remain differentiable
functions of actor parameters. Use reparameterized sampling.
The squashed-action log-density requires a change-of-variables correction:
\begin{equation}
\log\pi(a\mid s)=
\log\mathcal N(u;m_\theta,\operatorname{diag}(\sigma_\theta^2))
-\sum_j\log(1-\tanh^2(u_j)).
\end{equation}
Implement stable numerical forms and include any action-rescaling
Jacobian. Sum component log-densities to shape $[B,1]$.

\subsection{Critic Targets and Losses}
\begin{align}
a'&\sim\pi_\theta(\cdot\mid s'),\\
y&=r+\gamma(1-d)
[\min_iQ_{\bar\phi_i}(s',a')-\alpha\log\pi_\theta(a'\mid s')],\\
L_{Q_i}&=\E[(Q_{\phi_i}(s,a_{\rm stored})
-\operatorname{stopgrad}(y))^2].
\end{align}
The entire target computation, including current-policy sampling,
has no gradients. Slow target critics stabilize estimation.
Using the current actor reflects the latest policy, but its changes
and fresh samples can change labels even while target critics are fixed.
DDPG/TD3 choose additional slow actor copies to reduce action movement.
Neither choice is universally necessary; entropy alone does not
guarantee stability without a target actor.

With $r=2$, $\gamma=0.9$, $\alpha=0.2$, target Qs $(8,10)$,
and next log-density $-1$, $y=2+0.9(8+0.2)=9.38$.
Both online critics regress on the stored action toward this label.

\subsection{Actor Loss and Gradient Flow}
At replay states, sample fresh current actions with reparameterization:
\begin{equation}
L_\pi=\E_{s,\epsilon}
[\alpha\log\pi_\theta(a_\theta\mid s)
-\min_iQ_{\phi_i}(s,a_\theta)].
\end{equation}
Fix online critic parameters but preserve their action derivatives.
The negative Q-term favors valuable actions; expected log-density
encourages entropy. Standard SAC uses both online critics' minimum
in the actor objective; TD3 typically uses online Q1.
For online Qs $(7,9)$ and log-density $-0.7$, sample loss is
$0.2(-0.7)-7=-7.14$. Loss need not be nonnegative;
its derivatives determine the update.

PPO retains stored actions and recomputes their probabilities using
fixed rollout advantages. SAC samples new actions at replay states
for a different actor objective.

\begin{verbatim}
# actor.sample uses reparameterization and corrected log-density.
# actions: [B, act_dim]; log-density and Q: [B, 1]
with torch.no_grad():
    next_a, next_logp = actor.sample(next_states)
    next_q = torch.minimum(
        target_q1(next_states, next_a),
        target_q2(next_states, next_a))
    y = rewards + gamma * (1 - terminated.float()) * \
        (next_q - alpha * next_logp)
critic_loss = (
    (q1(states, stored_actions) - y).square()
    + (q2(states, stored_actions) - y).square()).mean()
# Clear gradients, backward, step critic optimizer.

# Freeze online critic parameters; retain action derivatives.
new_a, logp = actor.sample(states)
q = torch.minimum(q1(states, new_a), q2(states, new_a))
actor_loss = (alpha.detach() * logp - q).mean()
# Clear gradients, backward, step actor optimizer.
# Restore critic parameters; soft-update both target critics.
\end{verbatim}
Here alpha is a tensor; use a constant directly for fixed temperature.
Handle truncations separately from genuine terminal states and use
the correct final observation for bootstrapping.

\subsection{Automatic Temperature Adjustment}
A common extension learns alpha using a target average entropy
$\mathcal H_{\rm target}$. At fixed policy:
\begin{equation}
L_\alpha=\alpha(\E[-\log\pi]-\mathcal H_{\rm target}).
\end{equation}
Entropy below target increases alpha; entropy above target decreases it.
Optimize $\beta=\log\alpha$ so $\alpha=e^\beta>0$.
Writing $\hat{\mathcal H}=-\operatorname{stopgrad}(\log\pi)$,
the exact reparameterized dual loss is
$e^\beta(\hat{\mathcal H}-\mathcal H_{\rm target})$.
A common surrogate is
$\beta(\hat{\mathcal H}-\mathcal H_{\rm target})$:
same update direction but different gradient scaling.
\begin{verbatim}
alpha_loss = -(log_alpha *
    (logp.detach() + target_entropy)).mean()
# Step only the log_alpha optimizer.
\end{verbatim}
Temperature learning changes only the entropy weight.
Subsequent actor updates change its distribution; alpha does not
directly modify actor mean or variance.
After tanh, entropy is influenced by mean location as well as variance.
Target differential entropy can be negative.

\section{Comparison and Limitations}
\begin{center}
\small
\begin{tabular}{p{0.22\textwidth}p{0.20\textwidth}p{0.20\textwidth}p{0.22\textwidth}}
\toprule Feature & DDPG & TD3 & Modern SAC\\
\midrule
Actor & Deterministic & Deterministic & Stochastic\\
Online critics & One & Two & Two\\
Target actor & Yes & Yes & Usually no\\
Target critics & One & Two & Two\\
Actor objective & Maximize Q & Maximize Q1 & Maximize min Q + entropy\\
Exploration & Added noise & Added noise & Policy sampling\\
Special mechanisms & Slow targets & Twin Q, delay, smoothing & Entropy, optional learned alpha\\
\bottomrule
\end{tabular}
\normalsize
\end{center}
None guarantees accurate critics or global optimization.
Twin-Q minima may be pessimistic. Delay trades responsiveness for
stability. SAC still requires an entropy target when alpha is learned;
reward and action scaling influence its tradeoff.
A squashed Gaussian is a restricted policy family. Average entropy
does not guarantee useful exploration in every state or sparse-reward task.
Replay coverage, hyperparameters, and random seeds remain important.

\section{Additional Synthesis: Actions, Gradients, and Soft Bellman Learning}
\subsection{Three Actions, Three Different Roles}
For a replay transition $(s,a_{\mathrm{stored}},r,s',d)$, distinguish:
\begin{enumerate}
\item \textbf{Behavior action} $a_{\mathrm{stored}}$: the action that produced the observed reward and next state; critic regression evaluates this action.
\item \textbf{Target next action} $a'_{\mathrm{target}}$: generated at $s'$ by a target actor in DDPG/TD3 or sampled from the current actor in modern SAC, without gradients, to compute the TD label.
\item \textbf{Current actor action} $a_\theta(s)$: generated at replay state $s$ with an actor-gradient path; used in actor optimization.
\end{enumerate}
A stored action of $0.4$, a target next action of $0.62$, and a current actor output of $0.5$ can coexist in one training step. TD3 target smoothing noise is distinct from behavior exploration noise.

\subsection{Score Function versus Differentiating Through Actions}
In on-policy REINFORCE or PPO, the actor often uses
\begin{equation}
\nabla_\theta\log\pi_\theta(a_{\mathrm{stored}}\mid s)\hat A,
\end{equation}
where the action and advantage are fixed data. A deterministic DDPG actor instead receives gradients through the learned critic:
\begin{equation}
\nabla_\theta Q_\phi(s,\mu_\theta(s))=
\left.\nabla_a Q_\phi(s,a)\right|_{a=\mu_\theta(s)}
\nabla_\theta\mu_\theta(s).
\end{equation}
Modern SAC samples actions by reparameterization, so both the Q and log-density parts of its actor objective have gradients through actor parameters. Fix critic parameters during actor optimization but preserve its derivative with respect to action. Detaching the actor action would cut this path.
\begin{pitfall}
\texttt{critic.eval()} does not automatically disable critic parameter gradients. Freezing parameter gradients is different from detaching the critic output or actor action.
\end{pitfall}

\subsection{Deriving SAC's Soft Bellman Target}
The entropy-regularized state value for a fixed policy is
\begin{equation}
V^\pi_{\mathrm{soft}}(s)
=\E_{a\sim\pi_\theta(\cdot\mid s)}
[Q^\pi_{\mathrm{soft}}(s,a)-\alpha\log\pi_\theta(a\mid s)].
\end{equation}
Its soft Q-value obeys
\begin{equation}
Q^\pi_{\mathrm{soft}}(s,a)
=\E_{r,s'\mid s,a}[r+\gamma V^\pi_{\mathrm{soft}}(s')].
\end{equation}
Substitute the soft V definition, estimate the action expectation using a sample $a'\sim\pi_\theta(\cdot\mid s')$, and use target twin critics:
\begin{equation}
y_{\mathrm{SAC}}=r+\gamma(1-d)
[\min_iQ_{\bar\phi_i}(s',a')-\alpha\log\pi_\theta(a'\mid s')].
\end{equation}
Thus the log-density term is not an arbitrary adjustment: it follows from the entropy-regularized Bellman recursion. The modern variant evaluates this term with a current-policy action while target Q parameters change slowly.
For $r=1$, $\gamma=0.9$, target Q-values $(6,9)$, $\alpha=0.2$, and $\log\pi(a'|s')=-1.5$, the target is $1+0.9(6+0.3)=6.67$.

\subsection{Why Modern SAC Usually Has No Target Actor}
DDPG and TD3 use slow target actors to propose bootstrap actions. Modern SAC typically uses the current stochastic actor under no-gradient sampling with target critics. The critic label can therefore change from policy updates or resampling even if target critic parameters stay fixed. This is a design choice, not a theorem that entropy alone stabilizes targets. The actor still learns from current online critics and its entropy term; gradient signals are not supplied by a target actor.

\subsection{Unified Gradient and Target Checklist}
\begin{center}\small
\begin{tabularx}{\textwidth}{@{}lXXX@{}}\toprule
& DDPG & TD3 & Modern SAC\\\midrule
Actor type & Deterministic & Deterministic & Reparameterized stochastic\\
Critic regression action & Stored & Stored & Stored\\
Bootstrap action & Target actor & Smoothed target actor & Current actor, no grad\\
Target Q & One & Minimum of two & Minimum of two and entropy\\
Actor Q term & Online Q & Online Q1 & Minimum online Q\\
Actor other term & None & None & $\alpha\log\pi$\\\bottomrule
\end{tabularx}
\end{center}
The bootstrap expressions above all additionally include $r$, $\gamma$, and a terminal mask. In the critic loss, no gradient should flow through the bootstrap label. In the actor loss, gradient should flow through the current actor and its action into Q's action-input derivative.

\section{Review Questions}
\begin{enumerate}
\item Why is continuous-action maximization harder than DQN action selection?
\item Where does a learned critic's action-gradient signal come from?
\item How do freezing critic parameters and detaching actor actions differ?
\item Which actions enter critic regression, target construction, and actor training?
\item Why can a negative-Q actor loss coexist with a target actor?
\item How does Double DQN differ from TD3's clipped twin-Q target?
\item Why is delay about frequency rather than update ordering?
\item How do target smoothing and exploration noise differ?
\item Why can probability density exceed one and differential entropy be negative?
\item Why does SAC reparameterize actor-training samples?
\item Which computations use target critics and which use online critics?
\item Can SAC targets change while target critic parameters remain fixed?
\item What happens to alpha if policy entropy is below target?
\item Does an alpha update directly change the actor standard deviation?
\end{enumerate}

\section{Learning Status and Source Notes}
First study pass completed on October 13, 2026.
Based on the supplied continuous-action deep Q-functions material
and DDPG/TD3/SAC slides, supplemented by guided numerical examples
and code. These are study notes, not verbatim slide transcription
or coverage of every historical implementation.
Revisit tanh density corrections and temperature-loss conventions
during implementation practice.

\chapter{Model-Based RL: Dynamics, Planning, and Learning}
\section{Learning Goals and Method Landscape}
Model-based RL uses a model of the environment to help choose actions or improve learning. The CS 8803 Lectures 7--8 distinguish \emph{decision-time planning} from \emph{background planning}. We follow that framework:
\begin{itemize}
\item \textbf{Decision-time}: shooting, Model Predictive Control (MPC), Cross-Entropy Method (CEM), and probabilistic ensemble planning such as PETS.
\item \textbf{Background}: Dyna, short model-generated experience such as MBPO, model-based value expansion (MVE), and learning policies through differentiable models or planner supervision.
\end{itemize}
These labels are not mutually exclusive. A system can combine model-generated replay, multi-step value targets, and decision-time planning.
The primary lecture themes are dynamics models, planning with learned models, background versus decision-time planning, trajectory optimization, simulation for learning, terminal values, and policy backpropagation. PETS-specific TS1 and TS-infinity details below are supplementary paper context, rather than a verbatim treatment from the slides.

\section{What Does a Model Predict?}
A standard transition model represents
\begin{equation}
P(s_{t+1}\mid s_t,a_t),\quad
\hat s_{t+1}=f_\psi(s_t,a_t),\quad
\hat P_\psi(s_{t+1}\mid s_t,a_t).
\end{equation}
A model need not be a neural network: it can be known physics, a table, or a learned predictor. A learned deterministic model can use supervised regression:
\begin{equation}
L_{\mathrm{dyn}}(\psi)=\frac1N\sum_{i=1}^{N}
\|f_\psi(s_i,a_i)-s'_i\|_2^2.
\end{equation}
A probabilistic Gaussian model may predict the \emph{change} in state, $\Delta s=s'-s$, with mean $\mu_\psi(s,a)$ and diagonal covariance $\Sigma_\psi(s,a)$:
\begin{equation}
L_{\mathrm{NLL}}(\psi)=-\E_{(s,a,s')\sim\D_{\mathrm{real}}}
\log\hat P_\psi(s'-s\mid s,a).
\end{equation}
These model losses are learned from actual observed transitions. Dynamics models may predict rewards and terminations too, or those may be computed from known task functions. Model-free SAC learns policy and Q-values without needing an explicit transition predictor.

\section{Prediction Error and Compounding Error}
One-step accuracy does not guarantee reliable long-horizon predictions. Suppose uniformly
$\|\hat f(s,a)-f(s,a)\|\le\varepsilon$ and $f$ is $L$-Lipschitz in its state argument over the relevant region. For the same chosen action sequence, let $e_t=\|\hat s_t-s_t\|$. Then
\begin{align}
e_{t+1}&\le\varepsilon+Le_t,\\
e_H&\le\varepsilon\sum_{j=0}^{H-1}L^j\quad(e_0=0).
\end{align}
These are conditional bounds, not universal statements about learned neural dynamics. Error can feed future predictions inputs that were rare in the training set. An optimizer may also exploit systematic model errors: it maximizes $\hat J(A)$, not necessarily the true return $J(A)$. More thorough planning can raise predicted return and still lower real return.

\section{Decision-Time Planning: MPC versus CEM}
For a horizon $H$, plan an action sequence $A=(a_t,\ldots,a_{t+H-1})$ using the model:
\begin{align}
\hat s_t&=s_t,\quad
\hat s_{t+k+1}=f_\psi(\hat s_{t+k},a_{t+k}),\\
\hat J(s_t,A)&=\sum_{k=0}^{H-1}\gamma^k
\hat r(\hat s_{t+k},a_{t+k}),\\
A^*&\in\arg\max_A\hat J(s_t,A).
\end{align}
\textbf{MPC} is a receding-horizon \emph{control framework}: plan $H$ steps, execute only the first action, observe the next real state, and replan. \textbf{CEM} is one possible \emph{optimizer} inside the planner; MPC can instead use random shooting or gradient-based optimization, and CEM can optimize objectives outside MPC.

\subsection{Cross-Entropy Method Action-Sequence Search}
For $d_a$ action dimensions and horizon $H$, CEM fits a sampling distribution over $A\in\R^{H\times d_a}$. In the independent Gaussian variant, sample $N$ candidate action sequences from $q_\eta(A)$, score each through model rollouts, keep the top $K$ elites $\mathcal E$, and refit:
\begin{align}
\mu_{\mathrm{new}}&=\frac1K\sum_{j\in\mathcal E}A^{(j)},\\
\sigma_{\mathrm{new}}^2&=\frac1K\sum_{j\in\mathcal E}
(A^{(j)}-\mu_{\mathrm{new}})^2.
\end{align}
Clipping to action bounds, smoothing distribution updates, and imposing minimum variances are possible implementation choices. CEM is derivative-free and does not train the dynamics network during its iterations.
\begin{verbatim}
def cem_plan(s, dynamics, reward_model):
    mean, std = init_action_sequence_gaussian()
    for iteration in range(n_cem_iterations):
        sequences = sample_gaussian(mean, std, n_candidates)
        returns = simulate_and_score(s, sequences, dynamics, reward_model)
        elites = sequences[topk(returns, n_elite)]
        mean, std = elites.mean(axis=0), elites.std(axis=0)
    return mean[0]   # one possible final-action choice
\end{verbatim}
This represents the inner planner; the outer MPC loop executes one action and replans.

\subsection{CEM Numerical Check}
Let the two best action sequences be $(0.8,0.6)$ and $(0.4,1.0)$. The fitted mean is $(0.6,0.8)$, the elementwise maximum-likelihood variance is $(0.04,0.04)$, and the \emph{standard deviation} is $(0.2,0.2)$. The variance and standard deviation must not be confused. With 1000 candidates and 20 CEM iterations there are 20,000 candidate sequences sampled per planning call, under the assumed constant sampling count; MPC then executes one real action.

\section{Uncertainty and Probabilistic Dynamics Ensembles}
For multiple dynamics models $\hat f_{\psi_1},\ldots,\hat f_{\psi_M}$, differences in predicted conditional means can indicate \emph{epistemic} (knowledge) uncertainty. Predicted within-model stochastic variance can represent \emph{aleatoric} (irreducible outcome) uncertainty. In a Gaussian mixture approximation, total predictive variance separates into expected within-model variance plus variance of model means:
\begin{equation}
\mathrm{Var}(s'\mid s,a)\approx
\underbrace{\frac1M\sum_{i=1}^M\Sigma_i(s,a)}_{\text{within model}}
+
\underbrace{\frac1M\sum_{i=1}^M
(\mu_i(s,a)-\bar\mu)(\mu_i(s,a)-\bar\mu)^\top}_{\text{between models}}.
\end{equation}
This decomposition is an ensemble estimate, not a guarantee of perfectly calibrated uncertainty. All models can agree and still be wrong.

\subsection{Mean versus Conservative Planning}
For candidate sequence $A$, suppose the ensemble predicts returns $J_i(A)$. One score uses $\bar J(A)=\frac1M\sum_iJ_i(A)$. A conservative alternative is
\begin{equation}
J_{\mathrm{pess}}(A)=\bar J(A)-\beta\sigma_J(A),
\end{equation}
where $\sigma_J$ is the standard deviation of predicted returns. This is an \emph{optional} risk-sensitive heuristic, not an assertion that every ensemble method or PETS uses this penalty. Better optimization of an incorrect model can produce \emph{model exploitation}.

\section{PETS: Probabilistic Ensembles with Trajectory Sampling}
PETS combines probabilistic dynamics ensembles, sampling-based propagation of candidate trajectories, and MPC planning with an optimizer such as CEM. Each model can predict Gaussian distributions over next-state changes and is trained by likelihood on real transitions. For a fixed candidate action sequence, use multiple particles $p$:
\begin{equation}
s_{k+1}^{(p)}\sim
\hat P_{\psi_{b(p,k)}}(\cdot\mid s_k^{(p)},a_k).
\end{equation}
Score a candidate by average simulated return across particles. Model assignments can be resampled at each step (TS1), or fixed along an entire particle trajectory (TS-infinity). Fixing model identity does \emph{not} remove within-model transition randomness. These exact named trajectory-sampling schemes are supplementary PETS paper details beyond the lecture overview.
\begin{pitfall}
There are two distinct probability distributions: $\hat P_{\psi_i}(s'\mid s,a)$ is learned using real transitions and gradient descent; $q_\eta(A)$ is CEM's action-sequence sampling distribution, refit using elites. Running CEM does not automatically optimize $\psi_i$.
\end{pitfall}

\section{Background Planning: Dyna and Dyna-Q}
Dyna integrates real experience, model learning, and model-generated \emph{planning updates}. Classical tabular Dyna-Q keeps both a Q table and a model table $\mathcal M(s,a)$.
After a real transition $(s,a,r,s',d)$:
\begin{enumerate}
\item Use it to update $Q(s,a)$ with one-step Q-learning:
\begin{equation}
Q(s,a)\leftarrow Q(s,a)+\eta
[r+\gamma(1-d)\max_bQ(s',b)-Q(s,a)].
\end{equation}
\item Update the model table with the observed reward, successor, and termination information.
\item Repeat $n_{\mathrm{plan}}$ times: select a previously seen $(\tilde s,\tilde a)$, ask the model for $(\tilde r,\tilde s',\tilde d)$, and perform the \emph{same} Q-learning update on this simulated transition.
\end{enumerate}
In deterministic tabular Dyna, model storage can simply remember the latest observed transition. It need not use neural gradient descent. In that special setting model-based updates can coincide numerically with replaying historical transitions; more expressive stochastic or generalizing models make the distinction stronger.
\begin{verbatim}
for real_step in environment:
    a = epsilon_greedy(Q, s)
    s_next, r, terminated = env.step(a)
    q_update(Q, s, a, r, s_next, terminated)
    model[(s, a)] = (r, s_next, terminated)
    for _ in range(n_planning):
        s_sim, a_sim = sample_seen_pair(model)
        r_sim, next_sim, done_sim = model[(s_sim, a_sim)]
        q_update(Q, s_sim, a_sim, r_sim, next_sim, done_sim)
    s = s_next
\end{verbatim}
A planning update here means an extra simulated Bellman update, \emph{not} necessarily CEM action-sequence optimization. Simulated errors can bias values.

\subsection{Dyna Numerical Example}
Consider $s_0\xrightarrow{r=1}s_1\xrightarrow{r=5}\mathrm{terminal}$ with one action per state, $\gamma=0.9$, learning rate $1$, and zero initial Q. Two real interactions yield $Q(s_0)=1$ and $Q(s_1)=5$. A later model-based update of $(s_0,a_0)$ changes $Q(s_0)=1+0.9(5)=5.5$ \emph{without a third real interaction}. The benefit is extra computational learning from known experience, not newly observed reward.

\section{MBPO: Short Model Rollouts to Train SAC}
Model-Based Policy Optimization (MBPO) trains a dynamics ensemble on real data $\D_{\mathrm{real}}$, starts short model rollouts at states drawn from real replay, acts within the learned model using the current SAC policy, and stores simulated transitions in $\D_{\mathrm{model}}$. SAC minibatches mix real and model transitions.
\begin{align}
s_0&\sim\D_{\mathrm{real}},\quad
a_k\sim\pi_\theta(\cdot\mid s_k),\quad
\hat s_{k+1}\sim\hat P_\psi(\cdot\mid s_k,a_k),\\
\D_{\mathrm{SAC}}&=\text{minibatches from }
\D_{\mathrm{real}}\text{ and }\D_{\mathrm{model}}.
\end{align}
Using a mixed sample $(s,a,r,s',d)$, the ordinary one-step SAC target remains:
\begin{equation}
y_{\mathrm{SAC}}=
r+\gamma(1-d)\left[
\min_iQ_{\bar\phi_i}(s',a')
-\alpha\log\pi_\theta(a'\mid s')\right],
\quad a'\sim\pi_\theta(\cdot\mid s').
\end{equation}
Critic prediction uses the \emph{stored} current action $a$. Actor training samples a \emph{new} reparameterized action at the sampled replay state. Synthetic transitions are data, not action labels for imitation.
\begin{verbatim}
# Train dynamics on real data, not on its own synthetic labels.
dynamics_loss = nll(dynamics, real_transitions)
# Periodically generate short model rollouts, without gradients.
with torch.no_grad():
    model_transitions = rollout(actor, dynamics, real_start_states, k)
model_buffer.extend(model_transitions)
# SAC then uses mixed real+synthetic batches.
mixed = concat(real_buffer.sample(n_real),
               model_buffer.sample(n_model))
critic_step = standard_sac_critic_update(mixed)
actor_step  = standard_sac_actor_update(mixed.states)
\end{verbatim}
A short rollout reduces compounding model error; its optimal horizon is task- and model-dependent. More simulated data is not equivalent to more independent real information. The dynamics model is optimized by its own supervised loss, not normally by the SAC actor loss.

\section{MVE: Model-Based Value Expansion}
Model-Based Value Expansion (MVE) instead changes the \emph{value target}. Start with real $(s_t,a_t,r_t,s_{t+1})$, simulate additional rewards and transitions using the dynamics model and a chosen policy, then bootstrap at the predicted horizon:
\begin{equation}
y_t^{(H)}=
r_t+\sum_{k=1}^{H-1}\gamma^k\hat r_{t+k}
+\gamma^HV_{\mathrm{target}}(\hat s_{t+H}).
\end{equation}
A Q critic trains on the original $(s_t,a_t)$ via
\begin{equation}
L_Q(\phi)=\E[(Q_\phi(s_t,a_t)-\operatorname{stopgrad}(y_t^{(H)}))^2].
\end{equation}
For $\gamma=0.8$, rewards $(3,2,4)$ (first observed, later two modeled) and terminal predicted value $10$, the target is
$3+0.8(2)+0.8^2(4)+0.8^3(10)=12.28$.
If combined with SAC, the appropriate soft entropy terms must also be included consistently. MVE reduces dependence on an early value bootstrap but adds model error; it is not guaranteed to reduce total bias.
\begin{pitfall}
MVE and MBPO are different design choices: MVE changes how a label is constructed; MBPO changes which transitions are available for training. Both can coexist. In standard MVE critic regression, stop-gradient prevents the critic loss from training the dynamics model or the actor that generated the rollout.
\end{pitfall}

\section{Terminal Value Functions and Planner Distillation}
A finite-horizon planner can be myopic. Adding a learned terminal value estimates return beyond its search horizon:
\begin{equation}
\hat J_H(s,A)=
\sum_{k=0}^{H-1}\gamma^k\hat r(\hat s_k,a_k)
+\gamma^HV_\phi(\hat s_H).
\end{equation}
This expression resembles an MVE target but serves a different purpose: terminal value is used to \emph{rank action sequences}, while MVE uses a similar expression as a \emph{critic regression label}. With $H=3,\gamma=0.9$, immediate three-step score $2$ and terminal value $20$ give $2+0.9^3(20)=16.58$, better than a score $8$ with zero terminal value, if those predictions are reliable.

A planner can supervise a fast actor using $(s,a^*_{\mathrm{planner}})$ pairs:
\begin{equation}
L_{\mathrm{distill}}(\theta)=
\E_s\|\mu_\theta(s)-a^*_{\mathrm{planner}}(s)\|_2^2.
\end{equation}
This is a simple deterministic imitation-style example. Planning labels can be refreshed on states visited by the learned actor, connecting to dataset aggregation and DAgger. Planning may also be used for policy initialization, improvement, or target generation.

\section{Differentiable Dynamics and Policy Backpropagation}
When dynamics and rewards are differentiable, an actor can receive gradients directly through simulated transitions. For a fixed state $s$:
\begin{equation}
a=\mu_\theta(s),\qquad s'=f_\psi(s,a),\qquad
J(\theta)=R_\omega(s').
\end{equation}
By the chain rule:
\begin{equation}
\frac{dJ}{d\theta}
=
\frac{\partial R_\omega}{\partial s'}
\frac{\partial f_\psi}{\partial a}
\frac{\partial\mu_\theta}{\partial\theta}.
\end{equation}
For multi-step rollouts, derivatives also propagate through each intermediate state and actor application (backpropagation through time). Long Jacobian products can vanish or explode, and the gradient follows the learned model, which can be biased.
\begin{verbatim}
for p in dynamics.parameters():
    p.requires_grad_(False)  # freeze model weights
s = initial_state
return_pred = 0
for k in range(horizon):
    a = actor(s)
    next_s = dynamics(s, a)  # NO no_grad(), NO detach()
    return_pred += (gamma**k) * reward_fn(next_s)
    s = next_s
actor_loss = -return_pred.mean()
actor_optimizer.zero_grad()
actor_loss.backward()
actor_optimizer.step()
\end{verbatim}
Freezing dynamics parameters does \emph{not} stop gradients through its action input. By contrast, wrapping the dynamics call in \texttt{torch.no\_grad()} severs the actor's gradient path. This differs from MBPO's data generation and MVE's detached target computation.

\subsection{Policy Backprop Numerical Example}
Set $a=\theta s$, $s'=2s+a$, $R(s')=-(s'-10)^2$, $s=2$, $\theta=1$. Then $a=2$, $s'=6$ and
\begin{equation}
\frac{dJ}{d\theta}
=\underbrace{-2(6-10)}_{8}
\underbrace{\frac{\partial(2s+a)}{\partial a}}_{1}
\underbrace{\frac{\partial(\theta s)}{\partial\theta}}_{2}
=16.
\end{equation}
With gradient-ascent step size $0.1$, $\theta_{\mathrm{new}}=1+0.1(16)=2.6$. The model parameters can remain frozen throughout.

\section{Data, Targets, and Gradient Summary}
\begin{center}\small
\begin{tabularx}{\textwidth}{@{}lXX@{}}\toprule
Method & Main model use & Updated parameters\\\midrule
PETS + MPC/CEM & Evaluate candidate trajectories, choose first action & Dynamics by real-data NLL at model training time; CEM sampling moments by elites during planning\\
Dyna-Q & Generate extra transitions for Bellman updates & Q table by real and simulated TD; model table from observed transitions\\
MBPO + SAC & Generate short replay trajectories & Dynamics on real data; SAC actor/critics on mixed replay via their standard losses\\
MVE & Construct multi-step critic labels & Online critic for the detached value loss\\
Planner distillation & Provide supervised action labels & Student policy by supervised gradient descent\\
Policy backprop & Differentiate predicted reward through dynamics & Actor gradients through model inputs; model parameters can remain frozen\\\bottomrule
\end{tabularx}
\end{center}
For all implementations distinguish actual observed states from synthetic states, the online Q predictor from no-gradient target branches, and model fitting from planning optimization.

\section{Boundary Conditions and Limitations}
\begin{itemize}
\item A time-limit truncation is not necessarily a terminal MDP state; use the proper final observation and bootstrap mask.
\item Reward and termination mechanisms must be modeled or otherwise computed for synthetic transitions.
\item Planning and training can exploit shared systematic model errors even when ensemble variance is small.
\item Decision-time planning can be computationally intensive; distilled policies trade planning computation for additional training and approximation.
\item Sample efficiency does not imply low wall-clock computation or universal final-return superiority.
\end{itemize}

\section{Further Reading: Topics Named in Lectures 7--8}
The course planning landscape also names model-predictive path integral and other shooting techniques, gradient-based trajectory optimization, direct collocation, MCTS, implicit planning, stochastic value gradients, and Dreamer-style approaches. We have \emph{not} completed paper-level derivations or implementations of all these methods; they are useful topics for later expansion. Their appearance here does not mark them as mastered.

\section{Final Review Questions}
\begin{enumerate}
\item How is MPC different from CEM? Why does MPC execute only one of the optimized actions?
\item Which Gaussian in PETS is fitted with transition NLL, and which is refit with elite sequences?
\item How do TS1 and TS-infinity propagate model uncertainty differently?
\item Why can longer planning or model rollouts make performance worse?
\item Explain how Dyna-Q obtains a better Q estimate without additional real transitions.
\item Why is classical deterministic tabular Dyna sometimes equivalent to experience replay?
\item Compare MBPO and MVE by \emph{training data source} versus \emph{value target}.
\item Why does MBPO actor optimization use a newly sampled actor action rather than imitate a stored model-buffer action?
\item Explain the difference between Terminal Value Planning and MVE label construction.
\item In model-based policy backpropagation, why must the dynamics call remain differentiable even if model weights are frozen?
\item CEM elites $(0.8,0.6)$ and $(0.4,1.0)$: compute fitted mean, variance, and standard deviation.
\item For $a=\theta s$, $s'=2s+a$, $R(s')=-(s'-10)^2$, $s=2$, $\theta=1$, compute $dJ/d\theta$ and the new $\theta$ after an ascent step of $0.1$.
\end{enumerate}
\section{Study Status and Source Notes}
First guided pass and three-module review completed October 13, 2026: algorithm recognition, data/gradient flow, and numerical derivations. The review corrected a variance-versus-standard-deviation confusion and verified the later MVE target calculation and one-step chain rule. Based primarily on the supplied Georgia Tech CS 8803 Fall 2026 Lectures 7--8. Dyna-Q implementation details, PETS-specific TS1/TS-infinity, explicit SAC/MBPO tensor flow, and worked numbers are guided-study expansions; not a verbatim lecture transcription. For original-paper follow-up: Sutton (1990), Chua et al.\ (2018), Janner et al.\ (2019).

\chapter{Beyond Policy Gradients: First-Order RL and World Models}
\label{ch:beyond_pg}
\noindent\textit{Lecture-aligned study notes: CS 8803 Lecture 9 (Fall 2026).}

\section{Stochastic Optimization}
The full RL objective is $J(\theta)=\mathbb{E}_{\tau\sim p_\theta}[\sum_{t=1}^T\gamma^{t-1}r(s_t,a_t)]$. The lecture initially considers a one-action projectile problem. The flight may last two seconds, but the decision horizon is one. Contact with a wall makes the deterministic reward non-smooth. Randomizing the action smooths the expected objective.

For example, $r(a)=10\mathbf{1}\{a\geq5\}$, $a\sim\mathcal N(\theta,\sigma^2)$ gives
\[
J(\theta)=10\Phi((\theta-5)/\sigma),\qquad
J'(\theta)=10\varphi((\theta-5)/\sigma)/\sigma.
\]
The raw step function has no informative slope below five, but the randomized objective has a positive slope when $\sigma>0$. This is not a guarantee of global optimization.

\section{Score Function Versus First-Order Gradients}
REINFORCE follows the log-derivative identity:
\[
\nabla_\theta J=\mathbb E_{a\sim\pi_\theta}
[r(a)\nabla_\theta\log\pi_\theta(a)].
\]
It requires sampling rewards but no derivative of reward or environment. For a fixed-variance Gaussian $\mathcal N(\theta,\sigma^2)$ the sampled gradient is $r(a)(a-\theta)/\sigma^2$.

\paragraph{Multi-step distinction.}
The single-reward formula only applies to the one-step example. For the start-state discounted objective, the multi-step gradient uses future reward:
\[
\nabla_\theta J=\mathbb E\left[
\sum_{t=1}^{T}\nabla_\theta\log\pi_\theta(a_t|s_t)
\sum_{k=t}^{T}\gamma^{k-1}r_k\right].
\]
Equivalently the relative reward-to-go $G_t=\sum_{k=t}^{T}\gamma^{k-t}r_k$ has an outer discount factor $\gamma^{t-1}$. Baselines independent of sampled action can reduce variance.

\paragraph{Pathwise/FOG.}
Reparameterize as $a=g_\theta(\epsilon)$ with parameter-independent noise. If the reward-action path is differentiable,
\[
\nabla_\theta J=\mathbb E_\epsilon[
(\partial r/\partial a)(\partial g_\theta/\partial\theta)].
\]
For $a=\theta+\sigma\epsilon$ with fixed $\sigma$, $\partial a/\partial\theta=1$.

\paragraph{Illustrative comparison.}
Take $r(a)=-(a-5)^2$, $a\sim\mathcal N(\theta,1)$, $\theta=3$ and sample $a=4$. Then
\[
\hat g_{\mathrm{PG}}=-1,\quad
\hat g_{\mathrm{FOG}}=+2,\quad
\nabla_\theta J(3)=+4.
\]
They disagree on one sample yet agree in expectation. In this example, their variances are $74$ and $4$ respectively. This does not imply universally lower FOG variance, especially with contact and nonsmoothness.

\begin{pitfall}
The reward need not be differentiable for REINFORCE. The gradient is of the log policy, multiplied by the observed return. FOG instead differentiates a computational path from reward back to policy parameters.
\end{pitfall}

\section{BPTT Through Dynamics}
For $a_t=\pi_\theta(s_t)$, $s_{t+1}=f(s_t,a_t)$ and final reward $r(s_{t+1})$, a simple chain is
\[
\frac{dJ}{d\theta}
=\frac{\partial r}{\partial s_{t+1}}
\frac{\partial f}{\partial a_t}
\frac{\partial \pi_\theta}{\partial\theta}.
\]
For $s_1=0$, $a_1=\theta$, $s_2=2a_1$ and $r=-(s_2-10)^2$, at $\theta=3$ the gradient is $16$, encouraging increased action. Across a long horizon, BPTT multiplies dynamics Jacobians, giving vanishing/exploding gradients; contact makes first-order gradients particularly challenging. The lecture reports a Cartpole training-time illustration but not a universal runtime advantage.

\section{SHAC: Short-Horizon Actor-Critic}
SHAC (Xu et al., 2022, lecture slides 30--40) avoids differentiating a full, long trajectory. Use a short differentiable rollout and a terminal critic estimate:
\[
J_{\mathrm{SHAC}}(\theta)=\mathbb E[
\sum_{t=1}^{H}\gamma^{t-1}r(s_t,a_t)
+\gamma^H V_\psi(s_{H+1})].
\]
The slides use a shifted-index undiscounted equivalent. Actor parameters are updated through short simulated dynamics and terminal \emph{state}; critic weights remain fixed during actor backpropagation. Critic learning uses model-free TD($\lambda$), based on bootstrapped targets
\[
R_t^{(k)}=\sum_{j=0}^{k-1}\gamma^j r_{t+j}
+\gamma^k V_{\bar\psi}(s_{t+k}).
\]
Detaching the terminal value removes the long-horizon actor gradient, but \emph{does not} remove gradients from short-horizon rewards. SHAC mitigates long gradient chains but requires differentiable simulation and faces sim-to-real mismatch.

\section{World Models and PWM}
Instead of engineering a simulator, learn a differentiable surrogate in latent space:
\[
z_t=\operatorname{enc}_\phi(o_t),\quad
\hat z_{t+1}=f_\phi(z_t,a_t),\quad
\hat r_t=r_\phi(z_t,a_t),\quad
a_t\sim\pi_\theta(\cdot|z_t).
\]
The lecture uses an illustrative 512-dimensional latent. PWM (Policy Learning with Multi-task World Models) is presented as \emph{world model + SHAC}, with additional engineering. First learn a multi-task, multi-embodiment surrogate from offline data, then optimize per-task actors with short modeled rollouts and terminal values.

The lecture's schematic world-model objective is a weighted sum of latent transition error and reward prediction cross-entropy:
\[
\mathcal L_{\mathrm{WM}}=\mathbb E[
\sum_{t=0}^H\gamma^t(
\|\hat z_{t+1}-\operatorname{sg}(\operatorname{enc}(o_{t+1}))\|_2^2
+\operatorname{CE}(r_t,\hat r_t))].
\]
Here $\operatorname{sg}$ stops gradient; reward cross-entropy presumes a compatible prediction representation.

\paragraph{Smoothness is not fidelity.}
The lecture reports a case where optimization via a smooth surrogate performs better than through ground-truth dynamics. This is a specific empirical comparison, not universal superiority. A biased learned model can lead to a wrong action. If real $s'=s+a$, learned $\hat s'=s+2a$, and reward is $-(s'-10)^2$ at $s=0$, the model favors $a=5$ but reality favors $a=10$.

\section{Practical Comparison and Review}
\begin{center}\small
\begin{tabularx}{\textwidth}{@{}p{0.18\textwidth}X X@{}}\toprule
Method & Gradient source & Important drawback\\\midrule
REINFORCE & Return-weighted score & Sampling variance\\
FOG & Differentiable action-reward path & Requires usable derivatives\\
BPTT & Full rollout derivatives & Long-horizon instability and contact\\
SHAC & Short simulated path + critic & Simulator cost, sim-to-real gap\\
PWM & Short learned path + critic & Data/training cost, model bias\\\bottomrule
\end{tabularx}
\end{center}
If a fast trustworthy simulator exists, SHAC avoids world-model fitting; if simulation is challenging but abundant offline data are available, PWM may be appealing. Neither method is unconditionally cheaper or more stable.

\paragraph{Review questions.} Why does multi-step REINFORCE use return-to-go? Why can two unbiased estimators have opposite sample gradients? What actor gradient remains when SHAC terminal value is detached? Can a smooth world model produce a systematically wrong policy?

\paragraph{References and scope.} Source: CS 8803 Lecture 9 (Fall 2026). Lecture references include Mohamed et al., \emph{Monte Carlo Gradient Estimation in Machine Learning} (2020); Xu et al., \emph{Accelerated Policy Learning with Parallel Differentiable Simulation} (2022); Suh et al., \emph{Do Differentiable Simulators Give Better Policy Gradients?} (2022). First-pass lecture notes, not a paper reproduction.

\chapter{Exploration in Deep Reinforcement Learning}
\label{ch:exploration}
\noindent\textit{Lecture-aligned study notes: CS 8803 Lecture 10 (Fall 2026).}

\section{Why Exploration Matters}
An agent must trade off exploiting actions that currently look good against learning about uncertain actions. Sparse rewards and long action sequences make naively random actions inefficient. In a stochastic multi-armed bandit, the expected cumulative regret is
\[
\operatorname{Reg}(T)=\sum_{t=1}^{T}(\mu^\star-\mu_{a_t}),
\]
where $\mu^\star$ is the highest true arm mean. Exploration can incur short-term regret while improving later decisions. In sequential RL the key additional issue is \emph{deep exploration}: reaching a useful state may require a coherent sequence of actions before any external reward appears.

The lecture organizes strategies broadly into optimistic exploration, posterior sampling, and information gain. Several deep-RL methods produce novelty bonuses rather than rigorous upper confidence bounds; these classifications are conceptual, not mathematical equivalences.

\section{Bandit Foundations: UCB, Thompson Sampling, and Information Gain}
\paragraph{Optimism through UCB.}
A common UCB rule for bounded stochastic rewards is
\[
a_t=\arg\max_a\left[\hat\mu_t(a)+\sqrt{\frac{2\log t}{N_t(a)}}\right].
\]
The empirical mean exploits knowledge; the bonus prefers arms with few observations. For example, at $t=100$, an arm with $\hat\mu=0.70,N=80$ has score about $1.039$, while an arm with $\hat\mu=0.55,N=5$ has score about $1.907$. The low-count arm is explored despite its lower mean. These values are illustrative scores, not probabilities.

\paragraph{Thompson sampling.}
Maintain a posterior over unknown arm parameters, sample a possible set of values $\tilde\mu_a$ from that posterior, then select $\arg\max_a\tilde\mu_a$. For Bernoulli rewards, a Beta prior is conjugate: $\operatorname{Beta}(\alpha,\beta)$ becomes $\operatorname{Beta}(\alpha+n_{\rm success},\beta+n_{\rm failure})$. A lower-mean but uncertain arm can occasionally win under a posterior draw. Unlike UCB, Thompson sampling does not simply add a fixed bonus to an estimated mean.

\paragraph{Information gain and IDS.}
A candidate action can be evaluated for expected reduction in uncertainty about an unknown variable $\Theta$:
\[
I(\Theta;Y\mid D,a)=H(\Theta\mid D)
-\mathbb E_{Y\mid D,a}[H(\Theta\mid D,a,Y)].
\]
Equivalently, this equals the expected posterior-to-prior KL divergence. Information-directed sampling (IDS) balances expected regret against information gain, rather than blindly maximizing entropy reduction. Information unrelated to good decisions can be less valuable.

\paragraph{Small entropy example.}
Let an unknown Bernoulli success rate belong to $\{0.2,0.8\}$ with equal prior probability. Prior entropy is one bit. After one success the posterior is $(0.2,0.8)$; after one failure it is $(0.8,0.2)$. Both have binary entropy $H_2(0.2)\approx0.722$, giving expected information gain $1-0.722\approx0.278$ bits.

\section{Count-Based Exploration in Large State Spaces}
Simple count bonuses include
\[
r_{\rm int}(s)=\frac{\beta}{\sqrt{N(s)}}.
\]
An exact count is useful in small discrete spaces, but near-identical observations in continuous or pixel-based environments may each have count one. Thus direct tabulation fails to generalize across similar states.

\subsection{Density Models and Pseudo-Counts}
Bellemare et al.\ motivate a generalized count from how an online density model changes when it sees state $s$. Let $\rho_n(s)$ be its probability estimate before an update, and $\rho'_n(s)$ its probability after updating once on $s$. Define equivalent counts by
\[
\rho_n(s)=\frac{\hat N_n(s)}{\hat n},\qquad
\rho'_n(s)=\frac{\hat N_n(s)+1}{\hat n+1}.
\]
Solving for the pseudo-count gives
\[
\boxed{\hat N_n(s)=
\frac{\rho_n(s)[1-\rho'_n(s)]}{\rho'_n(s)-\rho_n(s)}}.
\]
For $\rho_n=0.10$ and $\rho'_n=0.20$, $\hat N_n=0.8$. For $\rho_n=0.10$ and $\rho'_n=0.101$, $\hat N_n=89.9$. A density estimate that changes little after observing a state corresponds to a larger effective count. The formula requires appropriate learning-positive behavior ($\rho'_n(s)>\rho_n(s)$) and compatible probability assumptions; arbitrary density updates may give invalid values.

An exploration bonus can replace $N(s)$ with a regularized $\hat N(s)$, e.g.\ $r_{\rm int}\propto1/\sqrt{\hat N(s)}$. The slides mention density estimators such as Context Tree Switching (CTS). Generating visually convincing samples is not required: the model is used for probability estimates and their learning response.

\subsection{Hash-Based Counting}
Tang et al., \emph{\#Exploration}, map high-dimensional states into compact codes $h(s)\in\{0,1\}^k$ and count $N(h(s))$. The bonus may be proportional to $1/\sqrt{N(h(s))}$. Compression permits related observations to share a count, but code collisions can merge states with important different consequences. Learned compression may help, but it must preserve decision-relevant structure.

\section{Novelty Estimation: EX2, RND, and Model Prediction Error}
\subsection{EX2: Exemplar Classifiers}
EX2 (Fu et al.) asks whether a new exemplar state $s^\star$ is distinguishable from past observations. Train a classifier that treats the exemplar as positive and the historical state distribution as negative. With balanced class priors and a point-mass positive distribution in a discrete idealization, its optimal output at the exemplar satisfies
\[
D^\star_{s^\star}(s^\star)=\frac{1}{1+p_D(s^\star)},
\qquad
p_D(s^\star)=\frac{1-D^\star_{s^\star}(s^\star)}{D^\star_{s^\star}(s^\star)}.
\]
This yields a classifier-based density or novelty signal. Without suitable regularization, a high-capacity classifier may memorize pixel-level details and mark familiar states as novel. An amortized exemplar model shares computation instead of fitting an entirely new classifier from scratch for every observation.

\subsection{Random Network Distillation (RND)}
RND (Burda et al.) uses a fixed randomly initialized target network $f_{\bar\phi}(s)$ and a trainable predictor $g_\theta(s)$. Train only the predictor using
\[
\mathcal L_{\rm RND}(\theta)=
\mathbb E_{s\sim D}\|g_\theta(s)-f_{\bar\phi}(s)\|_2^2,
\quad
r_{\rm int}(s)=\|g_\theta(s)-f_{\bar\phi}(s)\|_2^2.
\]
Repeated training exposure often reduces error on familiar states; unfamiliar states may retain larger error. This is \emph{not} a visitation count or calibrated Bayesian uncertainty. Generalization, observation noise, and learning dynamics matter. In an RL pipeline the predictor learns its auxiliary regression task while the policy is updated using external and intrinsic rewards (possibly with separate normalization or value estimation).

\subsection{Prediction Error from Learned Dynamics}
The lecture also describes exploration bonuses from dynamics prediction errors, such as comparing predicted and observed next-state features:
\[
\hat z_{t+1}=f_\phi(z_t,a_t),\qquad
r_{\rm int}\propto\|\hat z_{t+1}-z_{t+1}\|_2^2.
\]
This differs from RND: RND predicts a frozen random feature mapping of a state; a dynamics predictor predicts an actual next-state representation. High prediction error can be driven by irreducible environmental randomness (the noisy-TV issue) instead of knowledge worth acquiring. RND may face related novelty failures but is not itself a next-state dynamics predictor.

\section{Posterior Sampling in Deep RL: Bootstrapped DQN}
Exact Bayesian sampling of neural-network Q-functions is difficult. Bootstrapped DQN (Osband et al.) approximates diverse plausible Q-functions with $K$ output heads, possibly sharing a representation. Random initialization and bootstrap data masks encourage differences among heads. For head $k$, a masked TD loss has the schematic form
\[
\mathcal L_k=
\sum_i m_i^{(k)}\left[
Q_k(s_i,a_i)-
\bigl(r_i+\gamma(1-d_i)\max_{a'}Q_k^{-}(s'_i,a')\bigr)
\right]^2,\quad m_i^{(k)}\in\{0,1\}.
\]
A transition generated by one head may still train other heads, subject to their masks, because DQN is off-policy.

At the start of each episode sample $k\sim\operatorname{Uniform}(\{1,\ldots,K\})$ and keep it fixed during the episode:
\[
a_t=\arg\max_a Q_k(s_t,a).
\]
This makes exploration \emph{temporally coherent}. A long corridor may require many consistently directed moves; resampling an independent random action or head at every step can produce oscillations. Choosing one head does not guarantee a particular action sequence or that the head is accurate. Bootstrap heads only approximate posterior diversity; they are not samples from a guaranteed exact Bayesian posterior.

\section{Information Gain in Deep RL: VIME}
VIME (Houthooft et al.) learns a Bayesian dynamics model $p_\theta(s_{t+1}\mid s_t,a_t)$ and rewards transitions that change beliefs about its parameters $\theta$. Given prior data $D$ and observed transition $x=(s_t,a_t,s_{t+1})$, realized information gain is
\[
D_{\rm KL}\left(p(\theta\mid D,x)\,\|\,p(\theta\mid D)\right).
\]
Expected information gain averages this KL over the possible $s_{t+1}$. Since exact neural-network posteriors are intractable, VIME approximates them by variational distributions, such as a factorized Gaussian
\[
q_\phi(\theta)=\prod_{j=1}^{d}\mathcal N(\theta_j;\mu_j,\sigma_j^2),
\]
and uses an approximate intrinsic bonus
\[
r_{\rm int}(s_t,a_t,s_{t+1})=
\beta D_{\rm KL}(q_{\phi'}(\theta)\|q_\phi(\theta)).
\]
For an illustrative one-dimensional update $q_\phi=\mathcal N(0,1)$ to $q_{\phi'}=\mathcal N(1,1)$, the KL is $0.5$ nats. The model and variational posterior must themselves be fitted from transition data, making VIME more involved than an RND prediction loss.

\begin{pitfall}
\textbf{Prediction error is not information gain.} A known stochastic transition kernel may have large realized prediction error but almost no posterior update. In principle Bayesian dynamics information gain focuses on uncertainty reducible by data (epistemic uncertainty), whereas irreducible outcome noise (aleatoric uncertainty) need not be informative. The quality of this distinction depends on the dynamics and inference approximations.
\end{pitfall}

\section{Unified Comparison and Practical Questions}
\begin{center}\small
\begin{tabularx}{\textwidth}{@{}p{0.18\textwidth}X X@{}}\toprule
Method & Exploration signal / action rule & Learned component and caveat\\\midrule
UCB & Mean plus confidence bound & Arm counts and means; bound assumptions\\
Thompson sampling & Greedy under sampled posterior parameters & Posterior over rewards; sampling quality\\
Pseudo-count & Bonus based on equivalent density count & Density estimator; learning-positive updates\\
Hash count & Bonus based on code visits & Hash/compression; collisions\\
EX2 & Exemplar discriminability & Classifier; memorization\\
RND & Error predicting frozen random features & Predictor; error is a novelty heuristic\\
Bootstrapped DQN & One sampled Q-head for an episode & Multiple TD learners; approximate posterior\\
VIME & Dynamics posterior-change KL & Bayesian dynamics; variational approximation cost\\\bottomrule
\end{tabularx}
\end{center}

\paragraph{Intrinsic reward versus policy sampling.}
Pseudo-counts, hashes, EX2, RND, and VIME can supply bonuses in a shaped reward such as $r_{\rm total}=r_{\rm ext}+\beta r_{\rm int}$. Bootstrapped DQN instead changes the behavior policy by selecting a Q-head and need not change the environment reward. VIME's dynamics model is trained to estimate learning progress: using a model for exploration is not synonymous with planning through it using MPC or CEM.

\paragraph{Review prompts.}
Why is exact counting ineffective for continuous states? Derive the pseudo-count from before/after density estimates. When can a classifier incorrectly label familiar images as novel? Why is the RND target frozen? Why commit to one Q-head per episode? Explain a case where prediction error is high but VIME information gain is low.

\paragraph{References and scope.}
Primary source: CS 8803 Lecture 10, \emph{Exploration} (Fall 2026). Lecture-discussed work includes Bellemare et al., \emph{Unifying Count-Based Exploration and Intrinsic Motivation}; Tang et al., \emph{\#Exploration}; Fu et al., \emph{EX2}; Burda et al., \emph{Exploration by Random Network Distillation}; Osband et al., \emph{Deep Exploration via Bootstrapped DQN}; Houthooft et al., \emph{VIME}; and Stadie et al.\ on exploration via learned model error. These notes are an instructional synthesis, not a reproduction of paper experiments.

\chapter{Offline Reinforcement Learning}
\label{ch:offline_rl}
\noindent\textit{Lecture-aligned study notes: CS 8803 Lecture 11, Offline RL and (Deep) Offline RL (Fall 2026).}

\section{Problem Setup and the Main Difficulty}
Offline reinforcement learning seeks to learn a policy using only a fixed dataset
\[
\D=\{(s_i,a_i,r_i,s'_i,d_i)\}_{i=1}^{N},
\]
typically collected by a behavior policy $\pi_\beta$ or a mixture of behavior policies. During offline training, no new transitions can be requested. This is not synonymous with \emph{off-policy}: DQN and SAC can reuse data from a different policy while still collecting new transitions during online learning.

Behavior cloning (BC) maximizes $\E_{\D}[\log\pi_\theta(a\mid s)]$ without using rewards or Bellman backups. Offline RL additionally exploits observed rewards and long-horizon value estimation. It may \emph{stitch} useful portions of different trajectories where the state coverage permits, but it cannot reliably invent actions absent from the dataset.

\paragraph{Distribution shift and bootstrapping.}
A learned critic may assign arbitrarily optimistic values to out-of-distribution (OOD) state-action pairs. A greedy policy then exploits these errors:
\[
y=r+\gamma(1-d)\max_{a'}Q(s',a').
\]
If an unsupported action has estimated value $100$ while supported actions have values $5$ and $8$, the maximum selects $100$ even if its true value is small. Bellman backups propagate this error to earlier transitions. An online agent may correct the estimate with fresh experience; an offline learner cannot. Thus \emph{policy improvement must respect data support, or value estimation must be conservative enough to avoid exploiting uncertain extrapolations}.

\begin{pitfall}
Off-policy does not imply offline safety. Replay buffers and target networks stabilize learning but do not, by themselves, solve the OOD-action problem. An offline algorithm may also suffer state-distribution shift, approximation error, and insufficient coverage of promising actions.
\end{pitfall}

\section{Off-Policy Evaluation (OPE)}
OPE asks for the return $J(\pi)$ of a target policy $\pi$ using trajectories generated by behavior policy $\pi_\beta$, without executing $\pi$.

\subsection{Trajectory Importance Sampling}
Under appropriate support and trajectory-distribution assumptions, the importance ratio is
\[
\rho(\tau)=\prod_{t=0}^{T-1}
\frac{\pi(a_t\mid s_t)}{\pi_\beta(a_t\mid s_t)},
\qquad
\widehat J_{\rm IS}=\frac1N\sum_{i=1}^N \rho(\tau_i)R(\tau_i).
\]
The environment transition probabilities cancel in the likelihood ratio when target and behavior policies share the same environment and initial-state distribution. Support is essential: if $\pi(a\mid s)>0$ but $\pi_\beta(a\mid s)=0$ on a relevant state, the corresponding ratio cannot be identified from behavior data. Ratios multiply across time; for example, a ratio of $2$ in each of five steps yields $32$. This can lead to extreme variance and low effective sample size.

\subsection{Doubly Robust Evaluation}
In the contextual-bandit setting a standard doubly robust form is
\[
\widehat V_{\rm DR}(s,a,r)
=\widehat V(s)+
\frac{\pi(a\mid s)}{\pi_\beta(a\mid s)}
\bigl(r-\widehat Q(s,a)\bigr),
\quad
\widehat V(s)=\E_{a'\sim\pi(\cdot\mid s)}\widehat Q(s,a').
\]
The first term predicts return using a model; the second corrects its observed residual using importance weighting. Under suitable identification and estimation assumptions, correctness of either the propensity model or the outcome model can suffice for consistency. This does not guarantee low finite-sample variance. Sequential variants recursively combine value baselines and importance-weighted temporal corrections.

\subsection{Marginalized Importance Sampling}
Instead of multiplying policy ratios along entire trajectories, marginalized methods work with occupancy measures such as
\[
w(s,a)=\frac{d^\pi(s,a)}{d^{\pi_\beta}(s,a)}.
\]
Here $d^\pi$ is a normalized discounted occupancy measure when using discounted return conventions. The density ratio is estimated from distributional flow or stationarity relationships, as in the lecture's discussion of GenDICE. This can mitigate trajectory-product variance but estimating high-dimensional occupancy ratios remains difficult and requires coverage.

\section{Linear Offline Policy Evaluation: TD, LSTD, and LSTDQ}
\subsection{From a TD Target to an LSTD System}
For a fixed policy, approximate
\[
V_w(s)=\phi(s)^\top w,\qquad
\delta_i(w)=r_i+\gamma(1-d_i)V_w(s_i')-V_w(s_i).
\]
The one-step TD target is $y_i=r_i+\gamma(1-d_i)V_w(s_i')$ and the TD error is $\delta_i=y_i-V_w(s_i)$. For the exact value function the Bellman condition is
$\E[\delta\mid s]=0$, \emph{not} $\delta_i=0$ for every stochastic transition.

Linear semi-gradient TD(0) updates
\[
w\leftarrow w+\eta\,\delta_i(w)\phi(s_i),
\]
since $\nabla_w V_w(s_i)=\phi(s_i)$ and the bootstrap target is treated as fixed during the individual parameter update. LSTD instead solves the empirical feature-weighted equilibrium
\[
\sum_i\phi(s_i)\delta_i(w)=0,
\]
leading to the linear system
\[
\boxed{Aw=b},\qquad
A=\sum_i\phi(s_i)\bigl[\phi(s_i)-\gamma(1-d_i)\phi(s_i')\bigr]^\top,\quad
b=\sum_i\phi(s_i)r_i.
\]
Solve $w=A^{-1}b$ only if $A$ is nonsingular; in practice a linear solver or regularization is preferable. Multiplication by $\phi(s)$ enforces one balance equation per feature, not a separate zero-residual equation per transition.

\paragraph{Two-transition worked example.}
Use scalar features with $\gamma=0.5$ and two nonterminal samples:
\[
(\phi(s_1),\phi(s_1'),r_1)=(1,0,1),\qquad
(\phi(s_2),\phi(s_2'),r_2)=(2,1,2).
\]
Then
\[
A=1(1-0)+2(2-0.5)=4,\qquad b=1(1)+2(2)=5,
\quad w=1.25.
\]
The individual TD errors are $\delta_1=1-1.25=-0.25$ and
$\delta_2=2+0.5(1.25)-2(1.25)=0.125$. Neither is zero, but
$1(-0.25)+2(0.125)=0$.

\subsection{Projected Bellman Equation}
Let $T^\pi$ denote the policy Bellman operator. The exact fixed point satisfies $V^\pi=T^\pi V^\pi$. A linear feature space generally cannot represent this function exactly. Standard on-policy linear LSTD seeks the projected fixed point
\[
\boxed{V_w=\Pi T^\pi V_w},
\]
where $\Pi$ is projection onto the feature span under the relevant state weighting. This implies a feature-weighted Bellman-residual orthogonality condition. It differs from directly minimizing $\sum_i\delta_i(w)^2$: the full gradient of that objective differentiates through the next-state value, whereas semi-gradient TD and LSTD use current-state features in their normal equations. For off-policy or insufficiently covered data, additional assumptions and stability issues arise.

\subsection{LSTDQ and Least-Squares Policy Iteration (LSPI)}
For action values use $Q_w(s,a)=\phi(s,a)^\top w$. To evaluate a fixed policy $\pi$, set
\[
A_\pi=\sum_i\phi(s_i,a_i)
\left[\phi(s_i,a_i)-\gamma(1-d_i)\phi(s_i',\pi(s_i'))\right]^\top,\quad
b=\sum_i\phi(s_i,a_i)r_i.
\]
This is LSTDQ. LSPI alternates \emph{policy evaluation} using LSTDQ and \emph{greedy policy improvement}:
\[
w_k\leftarrow \operatorname{LSTDQ}(\D,\pi_k),\qquad
\pi_{k+1}(s)=\arg\max_a\phi(s,a)^\top w_k.
\]
The dataset stays fixed, but $A_{\pi_k}$ changes because the bootstrap action $\pi_k(s_i')$ changes. Greedy improvement can select unsupported actions and thereby reintroduce distribution shift.

\section{Deep Offline RL: Restrict the Policy or Correct the Critic}
\subsection{Explicit Policy Constraints and BEAR}
A generic constrained improvement step is
\[
\max_\pi \E_{s\sim\D,\ a\sim\pi(\cdot\mid s)}Q(s,a)
\quad\text{subject to}\quad
D(\pi(\cdot\mid s),\pi_\beta(\cdot\mid s))\leq\epsilon.
\]
BEAR (Bootstrapping Error Accumulation Reduction) emphasizes staying near the action \emph{support} of the behavior distribution, using a maximum mean discrepancy (MMD) based approximation. Matching support is not the same as matching all behavior probabilities: a new policy may put more mass on better supported actions. MMD is a surrogate, not a strict per-action support certificate. Behavior-model error and overly weak or strong constraints can respectively permit extrapolation or prevent improvement.

\subsection{Implicit Constraints: AWR and Advantage-Weighted BC}
KL-regularized improvement admits an exponential-reweighting form
\[
\pi^*(a\mid s)\propto\pi_\beta(a\mid s)
\exp\!\left(\frac{A(s,a)}{\lambda}\right).
\]
This motivates advantage-weighted regression (AWR):
\[
\boxed{\mathcal L_{\rm AWR}(\theta)=
-\E_{(s,a)\sim\D}
\left[\exp\!\left(\frac{\widehat A(s,a)}{\lambda}\right)
\log\pi_\theta(a\mid s)\right].}
\]
The actor learns via weighted log likelihood of \emph{dataset actions}, rather than directly differentiating $-Q(s,\pi_\theta(s))$. High-advantage examples have stronger influence. The critic estimation step may still query policy-generated OOD actions in its Bellman target, however. This limitation motivates IQL.

\subsection{Implicit Q-Learning (IQL)}
IQL learns a value function, a Q-function, and a policy. Its key idea is to perform critic backups without sampling a new action from the learned actor at $s'$.

\paragraph{Value loss: asymmetric expectile regression.}
Define the residual $u=Q_{\bar\phi}(s,a)-V_\psi(s)$ and
\[
\ell_2^\tau(u)=
\begin{cases}
\tau u^2,&u\geq0,\\
(1-\tau)u^2,&u<0.
\end{cases}
\qquad
\boxed{\mathcal L_V(\psi)=\E_{\D}[\ell_2^\tau(Q_{\bar\phi}(s,a)-V_\psi(s))].}
\]
For $\tau>0.5$, underestimating high Q-values is penalized more than overestimating low Q-values, so $V_\psi(s)$ moves toward the upper expectile of \emph{in-dataset} action values. With equally weighted Q-values $2$ and $8$ and $\tau=0.8$, the optimum between them solves
$0.2(V-2)+0.8(V-8)=0$, giving $V=6.8$ rather than the mean $5$ or maximum $8$.

\paragraph{Q loss: value-based backup.}
\[
\boxed{\mathcal L_Q(\phi)=
\E_{\D}\left[
\left(Q_\phi(s,a)-
\bigl[r+\gamma(1-d)V_\psi(s')\bigr]\right)^2
\right].}
\]
The target contains $V(s')$ rather than $\max_{a'}Q(s',a')$ or an expectation over actions sampled from the new policy. It does not directly query an OOD action value in the bootstrap step.

\paragraph{Policy loss: advantage-weighted extraction.}
\[
\boxed{\mathcal L_\pi(\theta)=
-\E_{\D}\left[
\exp\bigl(\beta[Q_\phi(s,a)-V_\psi(s)]\bigr)
\log\pi_\theta(a\mid s)\right].}
\]
The exponential weight is treated as constant during the actor update. $\tau$ controls how optimistic $V$ is within the observed-action distribution; $\beta$ controls the strength of advantage preference in the actor. They are distinct knobs. In implementations, advantage-weight clipping, target Q-networks and double critics are common stability details. IQL cannot reliably discover wholly unsupported high-value actions, and errors in in-dataset Q estimates still matter.

\subsection{Conservative Q-Learning (CQL)}
Rather than avoiding OOD Q queries, CQL makes critic values conservative. For a discrete action space, a common form is
\[
\boxed{
\mathcal L_{\rm CQL}=
\mathcal L_{\rm TD}
+\alpha\left[
\E_{s\sim\D}\log\sum_a\exp Q(s,a)
-\E_{(s,a)\sim\D}Q(s,a)
\right].
}
\]
The log-sum-exp term pressures large action values downward, while the data-action term relatively preserves observed actions. This discourages a learned policy from exploiting spuriously high Q-values. Continuous-action variants approximate the action expectation using samples, rather than enumerate every action. Too little regularization fails to address overestimation; too much can suppress genuinely useful actions. CQL does not guarantee every OOD action's Q is below every in-dataset action's Q.

\section{Comparison, Training Pitfalls, and Extensions}
\begin{center}\small
\begin{tabularx}{\textwidth}{@{}p{0.15\textwidth}X X@{}}\toprule
Method & Main intervention & Important tradeoff or failure mode\\\midrule
BEAR & Restrict actor near behavior action support (MMD surrogate) & Coverage estimation; too strong a constraint reduces improvement\\
AWR & Weighted behavior cloning from advantages & Critic backups may still evaluate policy-generated OOD actions\\
IQL & Expectile $V(s)$, $r+\gamma V(s')$ backup, weighted BC & Quality of dataset actions and Q estimates; expectile and weight tuning\\
CQL & Conservative penalty on learned Q-values & Excessive pessimism and costly action sampling\\\bottomrule
\end{tabularx}
\end{center}

\paragraph{Model-based and offline-to-online extensions.}
The lectures also discuss COMBO, which combines learned dynamics and conservative offline value learning, and offline pretraining followed by online fine-tuning (including AWAC/IQL). Model-generated rollouts provide extra coverage but introduce model bias; online fine-tuning changes the setting by allowing new interaction.

\paragraph{Training checklist.}
Before choosing a method, check behavior support and dataset quality; distinguish action support from distribution matching; monitor extrapolation and critic divergence; tune $\tau$, advantage-weight temperature, MMD strength or CQL $\alpha$ with care; and do not assume that a lower training TD loss implies reliable evaluation on unseen actions.

\section{Review Questions}
\begin{enumerate}
\item Explain why off-policy learning and offline RL are not equivalent.
\item Show how OOD overestimation enters a max-Q Bellman target and propagates through time.
\item Why can trajectory importance sampling have high variance? State the support requirement.
\item Explain the correction term in a doubly robust estimator and how occupancy ratios differ from trajectory products.
\item Derive $Aw=b$ from $\sum_i\phi(s_i)\delta_i=0$ and explain why each $\delta_i$ need not vanish.
\item Why can LSPI change its LSTDQ system without changing the dataset?
\item Contrast BEAR's support constraint, IQL's expectile backup, and CQL's critic penalty.
\item For Q-values $2$ and $8$, why does expectile $\tau=0.8$ produce $V=6.8$?
\item Which of IQL's three losses uses dataset actions in policy extraction, and what differentiates it from SAC's actor update?
\end{enumerate}

\paragraph{Source and scope.}
Primary sources: CS 8803 Fall 2026 Lecture 11 slides, \emph{Offline RL} and \emph{(Deep) Offline RL}. Lecture-discussed algorithms include LSTD/LSTDQ/LSPI, BEAR, advantage-weighted policy learning, IQL, CQL, COMBO, and OPE approaches such as importance sampling and doubly robust estimation. The worked examples and explanations are instructional syntheses, not reproduced paper results.

\chapter{Imitation Learning and Inverse RL}
\section{Why Do We Need It?}
\TODO{Add a simple explanation and example.}
\section{Models and Inputs/Outputs}
\TODO{Describe networks and action selection.}
\section{Data, Targets, Losses, and Gradients}
\TODO{Add equations and training pseudocode.}
\section{Links to Other Methods}
\TODO{Compare with earlier algorithms.}
\section{Review Questions}
\TODO{Add our questions after the study session.}
\chapter{Unified Comparison：算法间的共同结构}
\section{问题动机（Why）}
\TODO{系统讲解之后填入课程对应的动机、适用条件与局限。}
\section{模型结构（Architecture）}
\TODO{梳理输入输出、各网络参数与动作选择方式。}
\section{数据与优化（Data, Objective, Gradient）}
\TODO{结合原始讲义、伪代码与训练数据流填入，标记 stop-gradient 与 target updates。}
\section{与前一算法的联系（Connections）}
\TODO{对比设计差异并解释为什么需要这一变化。}
\section{理解检查与遗留问题}
\TODO{每次讨论后更新掌握程度、反例与易混淆点。}
\chapter{Research Perspectives 与论文笔记}
\section{问题动机（Why）}
\TODO{系统讲解之后填入课程对应的动机、适用条件与局限。}
\section{模型结构（Architecture）}
\TODO{梳理输入输出、各网络参数与动作选择方式。}
\section{数据与优化（Data, Objective, Gradient）}
\TODO{结合原始讲义、伪代码与训练数据流填入，标记 stop-gradient 与 target updates。}
\section{与前一算法的联系（Connections）}
\TODO{对比设计差异并解释为什么需要这一变化。}
\section{理解检查与遗留问题}
\TODO{每次讨论后更新掌握程度、反例与易混淆点。}
\appendix
\chapter{Update Log}
\begin{itemize}
\item 2026-10-12 (v0.1): Created the first outline and fundamentals draft.
\item 2026-10-12 (v0.2): Rewrote the document in plain English; removed Chinese font dependencies; improved the Overleaf import instructions.
\item 2026-10-12 (v0.3): Expanded Chapter 1 with policy evaluation and improvement, Bellman and TD learning, Q target networks, GAE, and PPO actor-critic training.
\item 2026-10-13 (v0.4): Completed the first-pass Chapter 2 notes: REINFORCE, reward-to-go, baselines, actor-critic, GAE, importance ratios and PPO clipping.
\item 2026-10-13 (v0.5): Completed first-pass Chapter 3: Q-learning, DQN, experience replay, target networks, Double DQN, gradient flow, limitations, and actor-critic comparisons.
\item 2026-10-13 (v0.6): Completed first-pass Chapter 4: DDPG, TD3, SAC, target actors, critic/action gradients, delayed updates, entropy densities, and automatic temperature adjustment.
\item 2026-10-13 (v0.7): Added Chapter 3 SARSA and replay-distribution notes; added Chapter 4 action-role, gradient-flow and soft Bellman derivations.
\item 2026-10-13 (v0.8): Completed Chapter 5 model-based RL study guide, including MPC/CEM, ensembles, PETS, Dyna, MBPO, MVE, terminal values, policy distillation, and differentiable policy gradients.
\item 2026-10-20 (v0.9): Added Lecture 9 study chapter on score-function versus first-order gradient estimators, BPTT, SHAC, and PWM.
\item 2026-10-29 (v1.0): Added Lecture 10 exploration notes: UCB, Thompson sampling, IDS, pseudo-counts, hashes, EX2, RND, Bootstrapped DQN, and VIME.
\item 2026-10-31 (v1.1): Added Lecture 11 offline RL: OPE (IS, DR, marginalized IS), LSTD/LSTDQ/LSPI, distribution shift, BEAR, AWR, IQL, CQL, and COMBO.
\end{itemize}
\chapter{Course Materials}
\begin{itemize}
\item Lecture 0: Introduction and Logistics.
\item Lecture 1: Supervised Learning and Imitation Learning.
\item Lecture 2: Policy Gradients.
\item Lecture 3: Actor-Critic Algorithms.
\item Lecture 4: Value-Based Methods.
\item Lecture 5: Deep RL with Q-Functions.
\item DDPG, TD3, and SAC lecture slides.
\item Lectures 7 and 8: Model-Based Reinforcement Learning.
\item Lecture 9: RL beyond Policy Gradients.
\item Lecture 10: Exploration.
\item Lecture 11: Offline RL and (Deep) Offline RL.
\end{itemize}
\end{document}